% a02_Kpos: proof of Proposition prop:k01:Kpos ((K-pos)_N at zero scalar source)
\providecommand{\moving}[1]{\textcolor{orange!80!black}{[#1]}}
\providecommand{\Kpos}{\ensuremath{(\mathrm{K}\text{-}\mathrm{pos})_{N}}}
\providecommand{\Zk}{\ensuremath{Z_{k}}}
\providecommand{\nbar}{\ensuremath{\bar n}}
\providecommand{\Dsp}{\ensuremath{D^{\mathrm{sp}}}}
\providecommand{\DEG}{\ensuremath{(\mathrm{DEG})}}
\providecommand{\DEGB}{\ensuremath{(\mathrm{DEG}^{B})}}
\providecommand{\lay}[3]{\ensuremath{\mathcal L(#1;#2,#3)}}
\providecommand{\DelG}{\ensuremath{\Delta G}}

\section{Proof of \texorpdfstring{\cref{prop:k01:Kpos}}{the one-step bound on the normalised action}}\label{app:a02_Kpos}

This appendix proves \cref{prop:k01:Kpos}. Throughout, $k$ is the spacing index, $j$ the number of renormalisation steps, $\eta=5^{j-k}$ the block scale, and the two members of a diagonal pair are the $k$-model after $j$ steps and the $(k+1)$-model after $j+1$ steps (\cref{ssec:k01:setting}); quantities of the second member carry a prime. The scalar source vanishes, $s=u$; by \cref{rem:k01:notused}(i) nothing below concerns the vacuum constant $v_{k,j}$, and ``field-free, source-free'' components are dropped wherever they occur.

\subsection{Units, operators, and the decomposition of the effective action}\label{ssec:a02:setting}

\paragraph{Units.} Fix a depth $0\le i\le j$ and work in \emph{unit units of depth $i$}: the lattice $X_i$ of the $i$-th block variables has spacing $1$, the physical length unit is $l_i\defeq5^i\eps_k$ ($l_j=\eta$), the fine lattice has spacing $5^{-i}$, $e_i\defeq el_i^{1/2}$, $m_i^2\defeq m^2l_i^2$, $\mu_i^2\defeq\lmu^2l_i^2$, and the rescaled fine field is $A_i(x)\defeq l_i^{1/2}A(l_ix)$, so that a unit charge sees $U(A_b)=\exp(e_i5^{-i}qA_{i,b})$, $q$ the real antisymmetric generator with $q^2=-1$. At depth $j$ we write $\rho_i\defeq5^{i-j}$ for the scale of the slice of index $i\le j$ ($\rho_j=1$). The $(k+1)$-model at depth $i+1$ lives on the same lattice $X_i$ with fine spacing $5^{-i-1}$. All statements hold at every point $(e,\lambda,m^2,\lmu)$ of King's range; constants depend on it, on $J$, on King's $b_0$, $p$ and on the order $\nbar=13$ of his expansion, and on nothing else unless said.

\paragraph{Operators} (\cite[(2.10)--(2.17)]{K1}, \cite[(3.2)]{K2}). On the cube $\CK$ the scalar has the Dirichlet condition of \cref{def:s02:Kcube} and the bonds to the Dirichlet planes are kept; $H_A\defeq D_A^*D_A$ is the covariant Laplacian, $M_0(A)\defeq H_A+m_i^2$. $Q_i(A)$ averages the fine field over the $i$-th blocks along King's hierarchical staircases $\Gamma^{(i)}_{y,x}$ (of length at most $3$ in unit units, inside the block) with weights $w(x)=5^{-3i}f(x)$, $f=1$ above the finest step and $f=f_D$ at the finest step, $\sum_xw=1$ and density at most $\varrho\defeq(5/4)^3$. Then
\[
M(A)\defeq M_0(A)+a_iQ_i(A)^*Q_i(A),\quad G_i(A)\defeq M(A)^{-1},\quad H_i(A)\defeq a_iG_iQ_i^*,\quad \Delta^{(i)}(A)\defeq a_i-a_iQ_iH_i(A),
\]
and with $Q$ the plain one-step average $X_i\to X_{i+1}$, $K_i(A)\defeq\Delta^{(i)}(A)+a5^{-2}Q(A)^*Q(A)=C^{(i)}(A)^{-1}$. The photon operators are the same with $f=1$, $m\to\lmu$, the Neumann condition of \cref{def:s02:Kcube}, and blocks that tile $\CJ$; the background is $A^{(j)}\defeq H^{\mathrm{ph}}_jA_j$. On the support of $\chi_j$ (King's \cite[(3.2)]{K1} with the face factor \cite[(5.41)]{K2}) $|A_j|,|dA_j|,|\phi_j|$ are bounded by powers of $p=p(\eta)$, and $\varpi_i\defeq e_i\sup|dA^{(j)}_i|$ is the smallness parameter.

\paragraph{The decomposition.} By \cite[(3.90)]{K1} the effective action after $j$ steps is
\begin{equation}\label{eq:a02:390}
S^{(j)}_k=\tfrac12\ip{A_j}{\Delta^{(j)}_{\mathrm{ph}}A_j}+\ln N_j+\tfrac12\ip{\phi_j}{\Delta^{(j)}(A^{(j)})\phi_j}+\ln N_j(0)-\ln\frac{Z_j(A^{(j)})}{Z_j(0)}+P^{(j),1},
\end{equation}
where $Z_j(A)$ is the normalised scalar Gaussian integral \cite[(2.15)]{K1} and $P^{(j),1}$ the sum of King's graphs \cite[Sect.~3.3]{K1} with the source vertex \cite[(3.20)]{K1} and the free-boundary mass counterterm of \cite[p.~330]{K2}. The constants $\ln N_j$, $\ln N_j(0)$ and the field- and source-free part of $P^{(j),1}$ belong to $v_{k,j}$. The remaining pieces, each compared across the diagonal pair on $\supp\chi_j$, are:
\begin{itemize}[nosep]
\item[(B1)] the photon quadratic form, $\frac12\ip{A_j}{(\Delta^{(j)}_{\mathrm{ph}}-\Delta'^{(j+1)}_{\mathrm{ph}})A_j}$;
\item[(B2)] the scalar quadratic form $\frac12\ip{\phi_j}{\Delta^{(j)}(A^{(j)})\phi_j}$;
\item[(B3)] the background determinant $-\ln[Z_j(A^{(j)})/Z_j(0)]$;
\item[(B4)] the graphs of $P^{(j),1}$ containing a curvature vertex;
\item[(B5)] the source-free, field-dependent graphs of $P^{(j),1}$ localised away from the faces (zone~I, \cref{ssec:a02:zones});
\item[(B5$'$)] the same, localised at the faces (zone~B);
\item[(B6)] the members of (B3)--(B5$'$) containing King's block vertices \cite[(3.25)--(3.26)]{K1}.
\end{itemize}
\Cref{prop:k01:Kpos} follows from \cref{prop:a02:B1,cor:a02:B2,cor:a02:B3,thm:a02:B4,thm:a02:B5,thm:a02:Bsc,thm:a02:BVP} and the position sum \cref{lem:a02:SC}, assembled in \cref{ssec:a02:assembly}.

\subsection{The photon quadratic form}\label{ssec:a02:B1}

Let $\Omega$ be a cube whose faces lie between lattice sites on block planes at every scale $\le j$ (this is $\CJ$: its side $2MJ$ is an integer), $2\Omega$ the doubled torus, $W$ the group generated by the reflections of $2\Omega$ in the faces of $\Omega$, and $E_f$, $E_c$ the even extensions of fine and coarse fields from $\Omega$ to $2\Omega$. King's photon action \cite[(2.5)]{K2} is componentwise, and the Neumann condition of \cref{def:s02:Kcube} keeps the gradient pairs inside $\Omega$: each component is a scalar field with the Neumann (induced-subgraph) Laplacian $M_\Omega$.

\begin{lemma}[even restriction]\label{lem:a02:P}
Let the photon blocks tile $\Omega$. Then $\Delta^{(i)}_{2\Omega}E_c=E_c\Delta^{(i)}_\Omega$, $G_{2\Omega}E_f=E_fG_\Omega$ and $H_{2\Omega}E_c=E_fH_\Omega$ (photon operators at $A=0$). Hence $\ip{A}{\Delta^{(i)}_\Omega A}_\Omega=2^{-3}\ip{E_cA}{\Delta^{(i)}_{2\Omega}E_cA}_{2\Omega}$, and every bound on the torus photon kernels of \cite[Thm.~3.3, Props.~3.8, 3.10]{K1} on $2\Omega$ holds verbatim for the Neumann photon on $\Omega$. Moreover $H_\Omega1=c_i1$ with $c_i=a_i/(a_i+\mu_i^2)$.
\end{lemma}
\begin{proof}
The Neumann Laplacian on $\Omega$ is the restriction of the torus Laplacian on $2\Omega$ to $W$-even fields (the reflection planes lie between sites), so $M_{2\Omega}E_f=E_fM_\Omega$. Since $W$ maps blocks to blocks, $Q_{2\Omega}E_f=E_cQ_\Omega$ and $Q_{2\Omega}^*E_c=E_fQ_\Omega^*$ (the adjoint of a $W$-covariant uniform block average is the uniform block extension, which maps even fields to even fields). Hence $(M_{2\Omega}+aQ^*Q)E_f=E_f(M_\Omega+aQ^*_\Omega Q_\Omega)$; inverting on the invariant even subspace gives $G_{2\Omega}E_f=E_fG_\Omega$, then $H_{2\Omega}E_c=aG_{2\Omega}Q^*E_c=aG_{2\Omega}E_fQ_\Omega^*=E_fH_\Omega$ and $\Delta_{2\Omega}E_c=aE_c-aQE_fH_\Omega=E_c(a-aQ_\Omega H_\Omega)$. The quadratic-form identity follows from $E_c^*E_c=2^3$ on even fields. Constants are Neumann-harmonic and $Q_\Omega1=1$, which gives $H_\Omega1=c_i1$. If the blocks do not tile $\Omega$ the identity fails at the faces; the tiling is part of the definition of the flow (\cref{ssec:k01:setting}), and the state does not depend on the flow.
\end{proof}

\begin{proposition}[(B1)]\label{prop:a02:B1}
On $\supp\chi_j$, $\big|\ip{A_j}{(\Delta^{(j)}_{\mathrm{ph}}-\Delta'^{(j+1)}_{\mathrm{ph}})A_j}\big|\le C5^{-2j}p^2(\lmu\eta)^{-2}\eta^{-3}|\CJ|$.
\end{proposition}
\begin{proof}
\cite[Prop.~3.10]{K1} bounds the difference of the Fourier symbols of $\Delta^{(j)}$ and $\Delta'^{(j+1)}$ on any torus that is a union of blocks at every scale by $C5^{-2j}$ relative to $\Delta^{(j)}(p)\le a_j$; its proof \cite[Sect.~4, Lemma 4.2]{K1} is the $f=1$ case of \cref{prop:a02:S}(a) below. Apply it on $2\CJ$, take the supremum over the momenta of $2\CJ$, and use \cref{lem:a02:P} with $\norm{E_cA_j}^2=2^3\norm{A_j}^2$ and $|A_j|\le Cp(\lmu\eta)^{-1}$ on $\supp\chi_j$ \cite[(3.2)]{K1}; the number of coarse sites is $\eta^{-3}|\CJ|$. This is \cite[(3.92)]{K1} with $|T|$ replaced by $|\CJ|$.
\end{proof}

\subsection{The scalar quadratic form}\label{ssec:a02:B2}

\begin{lemma}[the form is a two-leg tree]\label{lem:a02:Q2}
For every bond field $A$ and every $\phi_j$ on $X_j$,
\[
\ip{\phi_j}{\Delta^{(j)}(A)\phi_j}=a_j\norm{\phi_j}^2-a_j\ip{\phi_j}{Q_j(A)\phi^{(j)}},\qquad \phi^{(j)}\defeq H_j(A)\phi_j=a_jG_j(A)Q_j(A)^*\phi_j .
\]
\end{lemma}
\begin{proof}
$\Delta=a_j-a_j^2QGQ^*=a_j-a_jQ(a_jGQ^*)$.
\end{proof}
The second term is King's vertex \cite[(3.24)$\times$(3.25)]{K1} at $\phi'=0$ with no photon leg, and its external line is $\phi^{(j)}$, \cite[(3.28)]{K1}. The object has no internal line and no loop; its expansion in the background (\cref{prop:a02:DCONV}) consists of chains, every subgraph of which has non-negative degree, so it needs no counterterm and no face remainder. This is the content of \cite[p.~669]{K1}, ``a graph with two external scalar fields''.

\begin{lemma}[the photon background]\label{lem:a02:FA}
On $\supp\chi_j$: (i) $|A^{(j)}|_\infty\le C_Hp(\lmu\eta)^{-1}$; (ii) $|dA^{(j)}|_\infty\le C_Fp$ and $[dA^{(j)}]_\alpha\le C_Fp$, hence $|F(A^{(j)})|_\infty\le2C_Fp$; (iii) $A'^{(j+1)}$ and $A^{(j)}$ agree at the common fine sites up to $C5^{-\gamma j}p(\lmu\eta)^{-1}$ in sup norm and $C5^{-\gamma j}p$ in the first difference. In the units of depth $i\le j$, $\varpi_i\le(l_i/\eta)^2\varpi_j$ and $\varpi_j\le C_Fe\eta^{1/2}p$.
\end{lemma}
\begin{proof}
(i) is \cite[(3.7)]{K1} for $H^{\mathrm{ph}}$ with $f=Q^*A_j$. (ii): $H1=c_j1$ (\cref{lem:a02:P}), so $dHf(x)=dH(f-f(z_x))(x)$; the kernel of $dH$ decays like $Ce^{-\delta|x-z|}$ by \cite[(3.7)]{K1}, $|f(z)-f(z_x)|\le|dA_j|_\infty|z-z_x|_1$ with $|dA_j|\le p$ by \cite[(3.2)]{K1}; the H\"older part follows from \cite[(3.8)]{K1} in the same way. (iii) is \cite[(3.71)]{K1} applied to $A_j$ minus a constant, with the constant identity. On $\CJ$ use \cref{lem:a02:P}; across a reflection plane the even extension has zero difference, so $|dE_cA_j|=|dA_j|$. The scaling of $\varpi$ is the definition.
\end{proof}

\begin{lemma}[block coercivity]\label{lem:a02:COER}
Let $X$ be a lattice box with any Neumann or Dirichlet scalar condition, or a torus, tiled by unit blocks $\Delta$, with probability weights of density $\le\varrho$ and staircases of length $\le C_\Gamma$ inside the blocks. Put $C_P\defeq\frac14$, $c_0\defeq[\max\{(2+4(\varrho+1)^2)C_P,4/a\}]^{-1}$ and $\nu_0\defeq(c_0/(16(3+aC_\Gamma^2\varrho)))^{1/2}$. If $e_i|F(C)|_\infty\le\nu_0$, then $D_C^*D_C+aQ(C)^*Q(C)\ge c_0/4$; hence $M(C)\ge c_0/8$ whenever $|m_i^2|\le c_0/8$.
\end{lemma}
\begin{proof}
Drop the bonds between blocks and the Dirichlet boundary bonds (both non-negative): the form is $\ge\sum_\Delta q_\Delta$ with $q_\Delta(\psi)=\sum_{b\in\Delta}5^{-i}|U(C_b)\psi(b_+)-\psi(b_-)|^2+a|Q(C)\psi(y_\Delta)|^2$. \emph{Free case.} With $\bar\psi$ the plain block mean and $Q^w$ the weighted mean, $|Q^w\psi-\bar\psi|\le(\varrho+1)\norm{\psi-\bar\psi}$, and the discrete Neumann Poincar\'e inequality on a box of side $\le1$ has constant $\le C_P$ ($\lambda_2=4n^2\sin^2(\pi/2n)\ge4$); hence $\norm\psi^2\le(2+4(\varrho+1)^2)C_P\norm{d\psi}^2+4|Q^w\psi|^2$, i.e.\ $q^0_\Delta\ge c_0\norm\psi^2$. \emph{Gauge.} $q_\Delta$ is gauge invariant. In the comb gauge on the spanning tree of $\Delta$ rooted at $y_\Delta$ every non-tree bond carries the holonomy of its fundamental cycle, which bounds a surface of at most $2\cdot5^{i}$ plaquettes, so $|U(\tilde C_b)-1|\le2e_i5^{-i}|F|$, and a staircase picks up at most $2C_\Gamma e_i|F|$. Hence $\norm{D_C\psi}^2\ge\frac12\norm{d\tilde\psi}^2-12e_i^2|F|^2\norm\psi^2$ and $|Q\psi|^2\ge\frac12|Q^w\tilde\psi|^2-4C_\Gamma^2\varrho e_i^2|F|^2\norm\psi^2$, so $q_\Delta\ge(c_0/2-4(3+aC_\Gamma^2\varrho)e_i^2|F|^2)\norm\psi^2\ge\frac{c_0}4\norm\psi^2$.
\end{proof}

\begin{lemma}[one interpolation step]\label{lem:a02:V}
Let $C=B+A'$, $V\defeq M(B)-M(C)$, $\epsilon_1\defeq\sqrt3e_i|A'|_\infty$, $\epsilon_2\defeq C_\Gamma\varrho^{1/2}e_i|A'|_\infty$, $c'\defeq c_0/8$ and $D\defeq G(C)^{1/2}VG(C)^{1/2}$. Then $\norm D\le\theta\defeq2\sqrt2(\epsilon_1+\sqrt a\epsilon_2)c'^{-1/2}+(\epsilon_1^2+a\epsilon_2^2)c'^{-1}$.
\end{lemma}
\begin{proof}
With $E\defeq D_B-D_C$ and $F_Q\defeq Q(B)-Q(C)$: $V=D_C^*E+E^*D_C+E^*E+a(Q_C^*F_Q+F_Q^*Q_C+F_Q^*F_Q)$, $\norm E\le\epsilon_1$ (since $|1-U(A'_b)|\le e_i5^{-i}|A'_b|$) and $\norm{F_Q}\le\epsilon_2$ (staircase phases and Cauchy--Schwarz). For $\psi=G^{1/2}\varphi$, $\norm{D_C\psi}^2+a\norm{Q_C\psi}^2=\ip\psi{(M-m_i^2)\psi}\le2\norm\varphi^2$ and $\norm\psi^2\le\norm\varphi^2/c'$; insert into $|\ip\psi{V\psi}|$.
\end{proof}

\begin{proposition}[convergence of the scalar effective operator]\label{prop:a02:DCONV}
On $\supp\chi_j$, for $1\le i\le j$, with $\Delta^{(i)}_k(A^{(j)}_k)$ and $\Delta'^{(i+1)}_{k+1}(A'^{(j+1)}_{k+1})$ on the common lattice $X_i$, $\gamma'\defeq\min(\gamma,\alpha)$ and every $\varsigma>0$,
\[
\delta_\Delta(i)\defeq\norm{\Delta^{(i)}_k(A)-\Delta'^{(i+1)}_{k+1}(A')}_{\ell^2(X_i)}+\norm{\Delta^{(i)}_k(0)-\Delta'^{(i+1)}_{k+1}(0)}_{\ell^2(X_i)}\le C_\varsigma p^c(\lmu\eta)^{-1}(5^{-\gamma'i}+\eta^\varsigma).
\]
The same bound holds for the $s$-derivatives of order $\le2$ of $K_i(s\tilde A)$ in the gauge of \cref{lem:a02:PS}, with weights $e^{-\delta'|\cdot-x|}$.
\end{proposition}
\begin{proof}
\emph{Inputs.} The zero-background line bounds of the $f_D$ flow on $\CK$ and on boxes with faces on block planes: the slice bounds \cite[(3.63)--(3.65)]{K1} with decay and H\"older clauses, the slice convergence \cite[(3.73)--(3.75)]{K1} and the external-line convergence \cite[(3.71), (3.74)]{K1} at $n=1$. On the free lattice they are \cref{prop:a02:FL}; on the cube they follow by the odd image sum of \cref{prop:a02:S}(c) (the image terms obey the free bounds with decay in $|x-gy|$). \emph{Step 1 (Schur).} It suffices to show $|\delta\Delta(y,y')|\le C\theta_ie^{-\delta|y-y'|}$ with $\theta_i$ the right side. \emph{Step 2 (gauge).} Take $B\defeq A_k(x_0)$ at a common coarse site $x_0$ next to $y$, the same in both models, and $\lambda(z)\defeq B\cdot(z-x_0)$; on $\CJ$ this is a global gauge transformation, under which the Dirichlet condition and the weights are invariant. Then $\Delta(A)(y,y')=e^{eq\lambda(y)}\Delta(\tilde A)(y,y')e^{-eq\lambda(y')}$ with $\tilde A\defeq A-B$, $|\tilde A(z)|\le C_Fp|z-x_0|$ (\cref{lem:a02:FA}(ii), \cite[(3.45)]{K1}); the phases of the two models agree up to $e_i|B-B'||y-y'|\le C5^{-\gamma j}p^2(\lmu\eta)^{-1}$ by \cref{lem:a02:FA}(iii). \emph{Step 3 (expansion).} Taylor-expand $\Delta(s\tilde A)(y,y')$ at $s=0$ to order $n_T$, using the exact algebra \cite[(3.47)--(3.53)]{K1} at $B=0$. The coefficient of $s^m$ is a finite sum of zero-background chains $Q^{(e_0)}GV_1G\cdots V_rGQ^{(e_r)*}$, $V_i$ the vertices of \cite[(3.53)]{K1} with $\tilde A$-legs, $Q^{(e)}\in\{Q(0),F_2(\tilde A,0)\}$. The remainder has background lines $G(s\tilde A)$, whose Combes--Thomas decay follows from \cref{lem:a02:COER} ($M$ is finite-range plus a form-small perturbation) and whose vertex sizes follow from \cref{lem:a02:V}; it is $\le C^{n_T}(C\varpi_ip)^{n_T+1}e^{-\delta|y-y'|}\le C\eta^\varsigma$ for $n_T\ge2\varsigma+2$, separately in each model. \emph{Step 4 (chains).} Compare one chain across the two models at the common fine sites ($\eps_k\Z^3\subset\eps_{k+1}\Z^3$) with the discrepancy $d_1(u',u)\defeq\sup|u'-u|+\sup|d'u'-du|$, propagated along the chain: $G$ maps $|\cdot|_\infty$ to $C^{1,\alpha}$; a hopping vertex maps $C^{1,\alpha}$ to $C^\alpha$ with size $C\varpi_i$ and two-scale discrepancy $C\varpi_i[5^{-\alpha i}|u'|_{1,\alpha}+d_1(u',u)+d_1(\tilde A',\tilde A)]$ (because $|d^hv-d^{h/5}v|\le h^\alpha[dv]_\alpha$); lines are compared by \cite[(3.73)--(3.75)]{K1}; a $Q^*$-type output is always followed by a line, $GQ^*=a^{-1}H$, and \cite[(3.71)]{K1} applies, and for $F_2^*$ outputs the staircase-integrated lines \cite[(3.74)]{K1}; at the end vertex $Q(\tilde A)$ the staircases nest, $\Gamma'^{(i+1)}_{y,x'}=\Gamma^{(i)}_{y,x}\cup\Gamma_{x,x'}$ (the $(k+1)$-model's $i$-blocks are the $k$-model's $(i-1)$-blocks), so $|U(\tilde A'(\Gamma'))-U(\tilde A(\Gamma))|\le e_i(C_\Gamma\sup|\tilde A'-\tilde A|+3\cdot5^{-i}|\tilde A'|)$; the weights satisfy $\sum_{x'\in B^1(x)}w'(x')=5^{-3i}$ against $w(x)=5^{-3i}f_D(x)$, and the resulting $f_D$-weighted minus plain average over a finest block is $\le C5^{-\alpha i}$ since the integrand is Lipschitz there. After at most $n_T$ steps the chain discrepancy is $\le C_{n_T}(5^{-\gamma i}+5^{-\alpha i}+d_1(\tilde A',\tilde A))\varpi_i^rp^ce^{-\delta|y-y'|}$. No subgraph degree is used: $D^*$ always acts on a $C^{1,\alpha}$ function, so doubly differentiated lines never occur. \emph{Step 5.} Sum the finitely many chains and add Steps 2--3.
\end{proof}

\begin{corollary}[(B2)]\label{cor:a02:B2}
On $\supp\chi_j$, $|\mathrm{B2}-\mathrm{B2}'|\le\frac12\delta_\Delta(j)\norm{\phi_j}^2\le Cp^2\lambda^{-1/2}\eta^{-7/2}|\CJ|\,\delta_\Delta(j)$, with $\delta_\Delta(j)\le C_\varsigma p^c(\lmu\eta)^{-1}(5^{-\gamma'j}+\eta^\varsigma)$. In particular (B2) obeys the bound \eqref{eq:k01:Kpos}.
\end{corollary}
\begin{proof}
\Cref{prop:a02:DCONV} at $i=j$ and the bound $|\phi_j(y)|\le Cp[\lambda\eta]^{-1/4}$ of \cite[(3.2)]{K1}.
\end{proof}

\paragraph{The zero-background symbol.} The following statement, used above through the image sum, is the Dirichlet-weight form of \cite[Lemma~4.2, Prop.~3.10]{K1} and \cite[App.~A, (A.10)--(A.11)]{K2}.

\begin{proposition}[the free $f_D$ flow]\label{prop:a02:FL}
On $5^{-j}\Z^3$ at $A=0$, for $f\in\{f_D,1\}$, uniformly in $j\ge1$ and in the mass $m^2\eta^2\in[0,m^2]$: the weighted $\ell^1$ bounds of \cite[Thm.~3.3, (3.6)--(3.8)]{K1} for $G_j$, $DG_j$ and their H\"older quotients; the slice bounds \cite[Prop.~3.7, (3.62)--(3.65)]{K1}, including the slices of index $<0$ of the finer model; and the one-step comparison bounds \cite[Props.~3.8--3.9, (3.71), (3.73)--(3.75)]{K1} in all their lines.
\end{proposition}
\begin{proof}
This is King's Fourier argument \cite[Sect.~4]{K1} with the averaging weight $f$, as \cite[App.~A]{K2} prescribes. The symbol of $Q_j^*Q_j$ at a coarse momentum $z$ is $\sum_l|\hat f_j(z+l)|^2$ over the $5^{3j}$ aliases $l$, where $\hat f_j(q)=\prod_\mu\sigma_j(q_\mu)$ factorises, $\sigma_j$ being the product over the levels of the one-step symbols; for $f_D$ the finest factor has centre of mass at the geometric centre of the block, so $\hat f_j(q)=e^{-iq\cdot(\frac12,\frac12,\frac12)}R_j(q)$ with $R_j$ real and even. The alias majorant $|\hat f_j(z+l)|\le C\prod_\mu\min(1,|z_\mu|/|(z+l)_\mu|)$ (with the modulus $\frac14$ at the alias centres for $f_D$, in place of the $f=1$ value $0$) makes $\sum_l|\hat f_j(z+l)|^2\hat h(z+l)^{-1}$ converge uniformly, $\hat h$ the symbol of $-\Delta+m_j^2$, and $\Delta^{(j)}(z)=(a_j^{-1}+\sum_l|\hat f_j(z+l)|^2(\hat h(z+l)+a_j\sum_{l'}|\hat f_j(z+l')|^2)^{-1})^{-1}$ by \cite[(4.2)]{K1}; the symbols extend analytically to a strip $|\Im z|\le\kappa$ and are bounded below there, which gives the exponential decay of $G_j$, $\Delta^{(j)}$ and $C^{(j)}$ (a Combes--Thomas argument on the unit lattice for the operator form). The slices are the factorisation \cite[(2.17)]{K1}, $G_{(i)}=a_i^2\rho_i^{-4}(G_iQ_i^*)C^{(i)}(Q_iG_i)$, whose outer factors are the profiles $H_i$ bounded in $C^{1,\alpha}$ at scale $\rho_i$ by the same symbol bounds; this gives (3.62)--(3.65). The one-step comparison: $\hat f'_{j+1}=\hat f_j\cdot\sigma_{\text{fine}}$ and $a_{j+1}^{-1}=a^{-1}+a_j^{-1}5^{-2}$, so $|\Delta^{(j)}(z)-\Delta'^{(j+1)}(z)|\le C5^{-2j}\Delta^{(j)}(z)$ for real $z$ (the two symbols differ by one more alias level, whose relative size is $O(5^{-2j})$ because the centres of mass coincide for $f_D$ and move by half a fine spacing for $f=1$), and $|\Delta^{(j)}(y,y')-\Delta'^{(j+1)}(y,y')|\le C5^{-\gamma'j}e^{-\kappa|y-y'|}$ for $\gamma'<2$ by a contour shift; (3.71), (3.73)--(3.75) follow factor by factor from the slice factorisation, as in \cite[Sect.~4]{K1}.
\end{proof}

\begin{proposition}[the Dirichlet symbol on the cube]\label{prop:a02:S}
(a) $|\Delta^{(j)}(z)-\Delta'^{(j+1)}(z)|\le C5^{-2j}\Delta^{(j)}(z)$ for real $z$, hence $\norm{\Delta^{(j)}-\Delta'^{(j+1)}}_{\ell^2(\Z^3)}\le Ca5^{-2j}$. (b) $|\Delta^{(j)}(y,y')-\Delta'^{(j+1)}(y,y')|\le C5^{-\gamma'j}e^{-\kappa|y-y'|_1}$ for every $\gamma'<2$. (c) Let $\Omega$ be a box with faces on block planes whose Dirichlet planes carry the zeros of $f_D$ (in particular $\CK$), and $\Delta^{(j)}(\Omega)\defeq a_j-a_j^2Q_jG_j(\Omega)Q_j^*$ at $A=0$. Then $\Delta^{(j)}(\Omega)(y,y')=\sum_{g\in\cG}\operatorname{sgn}(g)\Delta^{(j)}(y,gy')$ over the group $\cG$ generated by the reflections in the Dirichlet planes, $\norm{\Delta^{(j)}(\Omega)-\Delta'^{(j+1)}(\Omega)}\le Ca5^{-2j}$, and the kernel bound (b) holds with $\kappa/2$. Consequently, on $\supp\chi_j$ the zero-background part of (B2) differs across the pair by at most $Ca5^{-2j}p^2\lambda^{-1/2}\eta^{-7/2}|\CJ|$.
\end{proposition}
\begin{proof}
(a), (b) are in the proof of \cref{prop:a02:FL}. (c): the Dirichlet Laplacian on $\Omega$ is the restriction of the free one to $\cG$-odd fields, and the $f_D$ blocks are mapped to blocks by $\cG$ (the Dirichlet planes carry the weight's zeros, which is what the choice of $f_D$ achieves \cite[(3.8)--(3.9)]{K2}), so $E_c\Delta^{(j)}(\Omega)=\Delta^{(j)}E_c$ with $E_c$ the coarse odd extension, exactly as in \cref{lem:a02:P} with signs; this is the image sum, absolutely convergent by (b). For the operator norm periodise to the doubled box: $\Delta^{(j)}(\Omega)$ is the restriction of the torus operator to odd fields, $\ip\psi{\Delta^{(j)}(\Omega)\psi}=2^{-3}\ip{E_c\psi}{\Delta^{(j)}_{2\Omega}E_c\psi}$ with $\norm{E_c\psi}^2=2^3\norm\psi^2$, and the difference has norm $\le\sup_z|\Delta^{(j)}(z)-\Delta'^{(j+1)}(z)|$. The last claim is \cref{cor:a02:B2}'s arithmetic with $\delta_\Delta$ replaced by $Ca5^{-2j}$.
\end{proof}

\subsection{The background determinant}\label{ssec:a02:B3}

Write $A\defeq A^{(j)}$ and, with $R\defeq1+\lfloor(\lmu\eta)^{-1}\rfloor$, $C_r\defeq(r/R)A$, $B_r\defeq C_{r-1}$, $V_r\defeq M(B_r)-M(C_r)$, $D_r\defeq G(C_r)^{1/2}V_rG(C_r)^{1/2}$ ($1\le r\le R$), King's interpolation \cite[(3.93)]{K1}. Determinants are real determinants on $\R^N$-valued fields.

\begin{lemma}[the identities \cite{K1}, (3.94)--(3.95)]\label{lem:a02:DI}
On every finite domain, $\ln[Z_j(A)/Z_j(0)]=\frac12\ln\det[G_j(A)G_j(0)^{-1}]=\frac12\sum_{r=1}^R\ln\det(I+D_r)$.
\end{lemma}
\begin{proof}
The normalising prefactor of \cite[(2.15)]{K1} does not depend on $A$, so the ratio of Gaussian integrals is $\det(G(A)/G(0))^{1/2}$; $G(C)G(B)^{-1}=I+G(C)V$ and $\det(I+XY)=\det(I+YX)$.
\end{proof}

\begin{lemma}[the interpolation steps are small]\label{lem:a02:D}
There is $\eta_0(e,\lmu,m,\lambda)$ such that for $\eta\le\eta_0$, on $\supp\chi_j$, $\norm{D_r}\le\theta\le Ce\eta^{1/2}p\le\frac12$ for every $r$. For symmetric $D$ with $\norm D\le\theta<1$ and $n\ge1$,
\[
\Big|\Tr\ln(I+D)-\sum_{q=1}^n\frac{(-1)^{q-1}}q\Tr D^q\Big|\le\frac{\Tr|D|^{n+1}}{(n+1)(1-\theta)}.
\]
\end{lemma}
\begin{proof}
$|A'|=|A|/R\le C_H5p$ by \cref{lem:a02:FA}(i) and $|F(C_r)|\le|F(A)|\le2C_Fp$ by (ii), so \cref{lem:a02:COER} applies for $\eta\le\eta_0$ and \cref{lem:a02:V} gives $\norm{D_r}\le\theta$. The series bound: diagonalise; for $|x|\le\theta$, $|\ln(1+x)-S_n(x)|\le\sum_{q>n}|x|^q/q$. (The one-sided inequality $\ln(I+D)\ge S_n(D)$ for even $n$ printed in \cite[(3.96)]{K1} fails when $D$ has a negative eigenvalue, e.g.\ $\ln\frac12<S_2(-\frac12)$; the two-sided form above is what is used.)
\end{proof}

\begin{lemma}[the remainder of the series]\label{lem:a02:REM}
Let $\mathsf q>3$ (King uses $\mathsf q=2\nbar+1=27$). Under \cref{lem:a02:D}, with $\epsilon_1,\epsilon_2$ as in \cref{lem:a02:V},
\[
\Tr|D_r|^{\mathsf q}\le C_{\mathsf q}\big[(\epsilon_1+\sqrt a\epsilon_2)(1+\epsilon_1+\sqrt a\epsilon_2)\big]^{\mathsf q}N|X_j|
\]
uniformly in the fine spacing, $|X_j|$ the number of unit sites; hence, for both members of the pair, $\sum_{r\le R}|\mathrm{Rem}_r|\le CR(Ce\eta^{1/2}p)^{27}N\eta^{-3}|\CJ|\le C\eta^{\sigma+1}|\CJ|$ for every $\sigma\le8$. This is the remainder line of \cite[(3.98)]{K1}, on the cube and on the torus.
\end{lemma}
\begin{proof}
(1) Schatten--H\"older, $\norm{XYZ}_{S_{\mathsf q}}\le\norm X\norm Y\norm Z_{S_{\mathsf q}}$ \cite[Sect.~VI.6]{RS1}, with the decomposition of $V$ in \cref{lem:a02:V} and $\norm{G^{1/2}D_C^*}\le\sqrt2$, $\norm{\sqrt aQ_CG^{1/2}}\le\sqrt2$, $\norm{G^{1/2}}\le c'^{-1/2}$, gives $\norm D_{S_{\mathsf q}}\le C(\epsilon_1+\sqrt a\epsilon_2)(1+\epsilon_1+\sqrt a\epsilon_2)\norm{G^{1/2}}_{S_{\mathsf q}}$ and $\norm{G^{1/2}}^{\mathsf q}_{S_{\mathsf q}}=\Tr G^{{\mathsf q}/2}$. (2) $M(C)\ge c'$ and $M(C)\ge H_C+m_i^2$ give $G\le2(H_C+c'/2)^{-1}$, so $\Tr G^s\le2^s\Tr(H_C+c'/2)^{-s}$ by operator monotonicity. (3) Diamagnetic bound: the Trotter product with entrywise domination (the diagonal parts of $H_C$ and $H_0$ agree, also with the Dirichlet bonds) gives $\norm{e^{-tH_C}(x,x)}\le e^{-tH_0}(x,x)$ \cite{BFS1}, hence $\Tr(H_C+c)^{-s}\le N\Tr(H_0+c)^{-s}$. (4) The eigenvalues of $H_0$ on a box of sides $l_\mu\ge1$ at spacing $h$ are $\ge\sum_\mu4n_\mu^2/l_\mu^2$, so $\Tr(H_0+c)^{-s}\le C_{s,c}\prod_\mu l_\mu=C|X_j|$ for $s>3/2$, uniformly in $h$. (5) Take $s={\mathsf q}/2$ and apply \cref{lem:a02:D} with $n=26$.
\end{proof}

\begin{remark}[the middle line of \cite{K1}, (3.98)]\label{rem:a02:FP}
The per-order bound printed there, $q^{-1}|\Tr D_r^q-\Tr D_r'^q|\le C(5^{-\gamma j}\eta^{-\beta}+\eta^{\sigma+1})|T|$, does not hold for $q=2$. For a constant $A_j$ on the torus (which lies in $\supp\chi_j$) and a fixed block scale, a Fourier computation on each eigenplane of $q$ gives $-\frac12\Tr D_r^2=-\frac12c_*\delta^2|X_j|5^{j}+O(|X_j|)$, $\delta=e_j|A'|$, $c_*=\frac4{(2\pi)^3}\int_{[-\pi,\pi]^3}\sin^2u_1(\sum_\mu4\sin^2\frac{u_\mu}2)^{-2}d^3u=0.1712\ldots$ (the resolvent identity $G=G_m-G_maQ^*QG$ shows that the block part, of rank $\le3N|X_j|$, contributes $O(|X_j|)$), so the difference across the pair grows like $5^{j}$; the $q=1$ term diverges with the opposite sign. Only the sum over $q$ converges. The argument below therefore does not expand in $q$ at a fixed background; it uses the scale decomposition of the determinant, in which each scale contributes $O(e_i^2|dA_i|^2)$ per site without any cancellation between orders.
\end{remark}

\begin{lemma}[two exact splittings]\label{lem:a02:SYL}
(i) If $M_0(A)>0$, then $\ln\det[G_j(A)G_j(0)^{-1}]=\ln\det_{X_j}[\Delta^{(j)}(A)\Delta^{(j)}(0)^{-1}]-[\ln\det M_0(A)-\ln\det M_0(0)]$. (ii) With $\Delta^{(0)}\defeq M_0$ and $C^{(i)}\defeq(\Delta^{(i)}+al_{i+1}^{-2}Q^*Q)^{-1}$ in physical units,
\[
\ln\det G_j(A)-\ln\det G_j(0)=\sum_{i=0}^{j-1}\big[\ln\det C^{(i)}(A)-\ln\det C^{(i)}(0)\big],
\]
and the Schur complement $al_{i+1}^{-2}-a^2l_{i+1}^{-4}QC^{(i)}Q^*$ equals King's $\Delta^{(i+1)}$ of \cite[(2.14)]{K1}, with $a_{i+1}^{-1}=a^{-1}+a_i^{-1}5^{-2}$.
\end{lemma}
\begin{proof}
(i) $\det M=\det M_0\det(I_c+aQM_0^{-1}Q^*)$, and Woodbury gives $(I_c+aQM_0^{-1}Q^*)^{-1}=I-aQGQ^*=\Delta/a$. (ii) Sylvester and Woodbury at each step give $\det C^{(i)}=\det\Delta^{(i+1)}_{\mathrm{rec}}(al_{i+1}^{-2})^{-\dim}/\det\Delta^{(i)}$, which telescopes to $\det\Delta^{(j)}_{\mathrm{rec}}/\det M_0$ up to an $A$-free constant; compare with (i). Integrating out $\phi_i$ in the composition of the Gaussian smearings with variances $a_i^{-1}l_i^2$ and $a^{-1}l_{i+1}^2$ gives the variance $a_{i+1}^{-1}l_{i+1}^2$, because $QQ^*=I$ for the plain one-step average and the averages compose, $QQ_i=Q_{i+1}$, the staircases being hierarchical; $f_D$ enters only $Q_1$, so the steps $i\ge1$ are plain \cite[(3.2)]{K2}.
\end{proof}

\begin{lemma}[coarse coercivity and decay at every scale]\label{lem:a02:CS}
In the units of depth $i$, if $e_i|F_i|\cdot25\le\nu_0$ and $|m_i^2|\le c_0/2$ (constants of \cref{lem:a02:COER} at block side $5$), then $c_K\le K_i(A)\le a_i+a5^{-2}$ with $c_K>0$ independent of $i$ and of the fine spacing. Moreover (Combes--Thomas) $\norm{e^wG_ie^{-w}}\le2/c'$ and $\norm{e^w(K_i+t)^{-1}e^{-w}}\le2/(c_K+t)$ for $|\nabla w|\le\delta_{CT}$, uniformly in $t\ge0$.
\end{lemma}
\begin{proof}
$\ip\psi{\Delta^{(i)}\psi}=\inf_\varphi[\ip\varphi{M_0\varphi}+a_i\norm{\psi-Q_i\varphi}^2]$; with the minimiser $\varphi$ and $r\defeq\psi-Q_i\varphi$, $\ip\psi{K_i\psi}=\ip\varphi{M_0\varphi}+a_i\norm r^2+a5^{-2}\norm{Q\psi}^2$. \Cref{lem:a02:COER} for the $(i+1)$-step average $Q_{i+1}=QQ_i$ at block side $5$ gives $c_0\norm\varphi^2\le\norm{D\varphi}^2+a'\norm{Q_{i+1}\varphi}^2$; with $\norm{Q_{i+1}\varphi}^2\le2\norm{Q\psi}^2+2\norm r^2$, $\norm{D\varphi}^2\le\ip\varphi{M_0\varphi}+|m_i^2|\norm\varphi^2$ and $\norm\psi^2\le2\varrho\norm\varphi^2+2\norm r^2$ this gives $\norm\psi^2\le C_K\ip\psi{K_i\psi}$; the upper bound is $\norm Q\le1$. Combes--Thomas: $M$ has range $\le1$, and $e^wMe^{-w}-M$ is form-small, $|\ip\psi{B\psi}|\le C\delta_{CT}(\norm\psi\norm{D\psi}+\norm\psi^2)$; combine with the lower bounds.
\end{proof}

\begin{lemma}[per-site bound in a local gauge]\label{lem:a02:PS}
For $x\in X_i$ let $g_x(A)\defeq\Tr_N[\ln K_i(A)](x,x)$, so that $\ln\det C^{(i)}=-\sum_xg_x$. Then $|g_x(A)-g_x(0)|\le CN(e_i|dA_i|_{\mathrm{loc},x})^2$, where $|\cdot|_{\mathrm{loc},x}$ is the supremum weighted by $e^{-\delta'|\cdot-x|}$.
\end{lemma}
\begin{proof}
(1) Gauge covariance of $K_i$ (all staircases transform at their endpoints) gives $g_x(A)=g_x(\tilde A_x)$ with $\tilde A_x\defeq A-A(x)$, a global gauge transformation on $\CJ$, and $|\tilde A_x(z)|\le|dA|_\infty|z-x|_1$. (2) Let $P$ be orthogonal with $PqP^{-1}=-q$. Then $K_i(-A)=PK_i(A)P^{-1}$, so $g_x(s\tilde A)$ is even in $s$ and $\partial_sg_x(0)=0$. (3) $g_x(\tilde A)-g_x(0)=\int_0^1(1-s)\partial_s^2g_x(s\tilde A)\,ds$ with $\partial_s^2\ln K=\int_0^\infty[RK''R-2RK'RK'R]\,dt$, $R=(K+t)^{-1}$. (4) With $w\defeq\delta_{CT}|\cdot-x|$, $|[RK''R](x,x)|\le\norm{e^wRe^{-w}}^2\norm{e^{-w}K''e^{-w}}$; $K'$, $K''$ are fine-level chains with lines $G_i(s\tilde A)$ (\cref{lem:a02:CS}) and vertices linear or quadratic in $\tilde A$, whose growing coefficients $|\tilde A_x(z)|\le|dA||z-x|$ are absorbed by $e^{-w}$ ($\sum_ze^{-2\delta_{CT}|z-x|}|z-x|^2\le C$), so $\norm{e^{-w}K''e^{-w}}+\norm{e^{-w}K'e^{-w}}^2\le C(e_i|dA|_{\mathrm{loc},x})^2$. (5) The $t$-integral converges like $\int(c_K+t)^{-2}dt$.
\end{proof}

\begin{theorem}[the background determinant converges]\label{thm:a02:DETC}
On $\supp\chi_j$, with $F_{\mathrm{ph}}\defeq\sup|dA^{(j)}|$ in physical units ($\le Cp\eta^{-3/2}$ by \cref{lem:a02:FA}):
\begin{itemize}[nosep]
\item[(i)] $|\ln[Z_j(A)/Z_j(0)]|\le CNe^2F_{\mathrm{ph}}^2|\CJ|\sum_{i<j}l_i\le CNe^2p^2\eta^{-2}|\CJ|$;
\item[(ii)] $\big|\ln[Z_j(A)/Z_j(0)]-\ln[Z'_{j+1}(A')/Z'_{j+1}(0)]\big|\le CNe^2F_{\mathrm{ph}}^2|\CJ|\big[\sum_{i=1}^{j-1}l_i\delta_{\mathrm{conv}}(i)+3\eps_k+\eta^\varsigma\big]$ with $\delta_{\mathrm{conv}}(i)\le Cp^c(\lmu\eta)^{-1}(5^{-\gamma'i}+\eta^\varsigma)$ from \cref{prop:a02:DCONV}; hence the left side is $\le CNe^2p^{c+2}(\lmu\eta)^{-1}\eta^{-2}(5^{-\gamma'j}+5^{-j}+\eta^\varsigma)|\CJ|$.
\end{itemize}
\end{theorem}
\begin{proof}
(1) \Cref{lem:a02:SYL}(ii) writes both members as sums over scales. (2) Pair the $k$-model's step $i$ with the $(k+1)$-model's step $i+1$ (same physical scale, same coarse lattice $X_i$); for $i\ge1$ both one-step averages are plain. (3) The unmatched steps, the $k$-model's step $0$ (weight $f_D$) against the $(k+1)$-model's step $1$ (plain), and the $(k+1)$-model's finest step $0$, are bounded by \cref{lem:a02:PS} alone, each by $CNe^2F_{\mathrm{ph}}^2|\CJ|\eps_k$. (4) \Cref{lem:a02:PS} per site, times $\#X_i=|\CJ|l_i^{-3}$ and $(e_i|dA_i|)^2=e^2l_i^4F_{\mathrm{ph}}^2$, gives (i). (5) For (ii) compare per site: with the expansion of \cref{prop:a02:DCONV}, Step 3, in $s$ (order $n_T$, remainder bounded separately), the difference of $\partial_s^2g_x$ between the two members is $\le CN(e_i|dA_i|)^2\delta_{\mathrm{conv}}(i)$; integrate in $s$ and sum over $x$ and $i$, using $\sum_il_i5^{-\gamma'i}\le C\eta5^{-\gamma'j}$.
\end{proof}
No Ward identity, renormalisation or expansion in orders enters; \cref{lem:a02:PS} is uniform up to the Dirichlet faces, since the Dirichlet condition only adds non-negative bonds, so no face remainder arises for (B3).

\begin{corollary}[(B3)]\label{cor:a02:B3}
On $\supp\chi_j$, $|\mathrm{B3}-\mathrm{B3}'|$ obeys the bound of \cref{thm:a02:DETC}(ii), which is of the form \eqref{eq:k01:Kpos}.
\end{corollary}

\subsection{The graphs: localisation, zones, and the reduction away from the faces}\label{ssec:a02:zones}

From here on everything is in the unit units of depth $j$: the live sites are $T_\eta\defeq X_j$, the fine spacing is $5^{-j}$, $F$ denotes the scalar Dirichlet planes of $\partial T_\eta$, and $t(x)\defeq\dist(x,F)\ge5^{-j}$ for a live site. King localises every graph of $P^{(j),1}$ \cite[p.~661]{K1}: $\square$ is the smallest union of large blocks (side $M$) containing its vertices, $\square''\subset\square'$ the smallest boxes containing the large blocks within $p$, resp.\ $2p$, of $\square$, truncated at $\partial T_\eta$; $\Omega\defeq\square'\cap T_\eta$ is a rectangular box with faces on block planes at every level (this is where the integer side $2MJ$ enters), and the internal lines are the slices $G_{(i)}(\square',0)$ of the $f_D$ flow with the Dirichlet condition on all of $\partial\Omega$, the photon slices with the Neumann condition, and the lines $\delta G_j(\square',A^{(j)})$, external in King's sense. A localisation is in \emph{zone~I} if $\square'\cap F=\emptyset$ and in \emph{zone~B} otherwise; the faces of $\Omega$ not on $F$ lie at distance $\ge2p$ from $\square$. Lines obey King's slice bounds \cite[(3.63)--(3.65)]{K1} on $\Omega$ up to and including the faces (\cref{prop:a02:S}(c) for the scalar, \cref{lem:a02:P} for the photon, both image sums of free kernels), and the external objects obey \cite[(5.2)--(5.3)]{K2}: this is King's Prop.~5.1 on boxes whose faces lie on block planes, which follows from the image sums and a bootstrap on the sup of $|u|+|D_Au|$ with the commutator localisation of \cite[Lemma~5.3]{K2}, with no a priori sup bound. The range of $j$ is $0\le j\le k-J_0$, $5^{-J_0}\le\delta_*$, with $\delta_*^{-1}>3p(\delta_*)+4M+1$.

\begin{lemma}[zones]\label{lem:a02:Z}
(i) Every localisation with $R\defeq\dist(\{y\},\{z\})\le p$ that contains a source vertex lies in zone~I. (ii) Every zone-B localisation containing a source vertex has $R\ge\eta^{-1}-2p-4M\ge\eta^{-1}/2$. (iii) Every source vertex $s$ has $t(s)\ge\eta^{-1}/2\ge1$ and $\dist(s,\partial\square'\setminus F)\ge2p-\sqrt3$.
\end{lemma}
\begin{proof}
The sources are supported at continuum $\ell^\infty$-distance $>1$ from $\R^3\setminus\CJ$, i.e.\ $\ge\eta^{-1}$ unit lengths from $\partial T_\eta$; after the summation by parts of \cref{lem:a02:SBP} the summed vertex sits on bonds at distance $\le1-\eps$ from nothing (its weight vanishes within $1-\eps$ of the faces). In (i) every vertex is within $p$ of the source vertex, $\square$ within $p+M$, $\square'$ within $3p+3M<\eta^{-1}$. (ii), (iii) are the same arithmetic.
\end{proof}

\begin{lemma}[zone~I is the free $f_D$ flow]\label{lem:a02:I}
Let $\tilde H$ be a graph localised in zone~I and $\tilde H^{\mathrm{free}}$ the same graph with every internal scalar slice replaced by the free $f_D$ slice on $5^{-j}\Z^3$ and every internal photon slice by its free slice. Then $E(\tilde H)=E(\tilde H^{\mathrm{free}})+\mathcal R$, where every term of $\mathcal R$ carries one image line of size $\le C\rho_i^{-1}e^{-4\delta_0p/\rho_i}$ (with derivatives and H\"older quotients), so that an image line has any prescribed degree and $|\mathcal R|\le Ce^{-\delta_0p/2}\times$(the bound of $\tilde H^{\mathrm{free}}$ with that line's degree raised). Since $p(\eta)=b_0(1+\ln\eta^{-1})^p$ with $p>1$, $e^{-\delta_0p/2}\le C\eta^{3+\sigma+1/2}$. The free-boundary counterterm of \cite[p.~330]{K2} renormalises $\tilde H^{\mathrm{free}}$ exactly, because the full free covariance does not depend on the averaging weight.
\end{lemma}
\begin{proof}
$G_{(i)}(\square',0)-G_{(i)}(0)$ is the signed sum over the non-identity elements $g$ of the reflection group of $\Omega$ of free slices $G_{(i)}(x,gy)$ (\cref{prop:a02:S}(c)); for $x,y\in\square$ and $\dist(\square,\partial\square')\ge2p$, $|x-gy|\ge4p$, and the free slice bounds give the size. The photon slices are the even image sums of \cref{lem:a02:P}, and in zone~I $\dist(\square,\partial T_\eta)\ge2p$ as well. The lines $\delta G_j(\square',A)$ carry $e^{-\delta_0\dist(x,\partial\square')}$ by \cite[(5.3)]{K2}, also $O(e^{-\delta_0p})$. Expanding each line as free plus image gives the decomposition.
\end{proof}

\paragraph{The curvature vertex.} Let $f$ be a real $2$-form of the family, $g^\eps$ its centre sampling, and $F(p)$ the plaquette curvature.

\begin{lemma}[the summed curvature vertex]\label{lem:a02:SBP}
(a) For every bond field $A$ on $\CJ$, $\sum_pg^\eps(p)F(p)=\sum_bw(b)A_b$ exactly, with $w(y,\mu)=\sum_{\nu>\mu}\eps^{-1}[g(p_{\mu\nu}(y))-g(p_{\mu\nu}(y-\eps e_\nu))]+\sum_{\nu<\mu}\eps^{-1}[g(p_{\nu\mu}(y-\eps e_\nu))-g(p_{\nu\mu}(y))]$, no boundary term, and $w=0$ on bonds within $1-\eps$ of $\partial\CJ$. (b) $d^*w=0$, and $\sum_bw(b)B_b=0$ for constant $B$. (c) $w(b)=W^\eps_\mu(m_b)$ with $W^\eps_\mu(x)=\sum_{\nu\ne\mu}\eps^{-1}[f_{\mu\nu}(x+\frac\eps2e_\nu)-f_{\mu\nu}(x-\frac\eps2e_\nu)]$, the restriction to bond midpoints of a smooth function with $\norm{\partial^\alpha W^\eps}_\infty\le2\max\norm{f_{\mu\nu}}_{C^{|\alpha|+1}}$ uniformly in $\eps$; so $w$ has sup and lattice $(1,\alpha)$-norms bounded by $C\max\norm{f_{\mu\nu}}_{C^{2,\alpha}}$. (d) For bonds $b'$ at spacing $\eps_{k+1}$ and $b$ at spacing $\eps_k$ in the same direction with $|m_{b'}-m_b|\le\sqrt3\eps_k$, $|w_{\eps_{k+1}}(b')-w_{\eps_k}(b)|\le(2\sqrt3+1)\eps_k\max\norm{f_{\mu\nu}}_{C^2}$.
\end{lemma}
\begin{proof}
(a) is a finite rearrangement: each $A_b$ collects the plaquettes containing $b$ with their orientation signs; a plaquette containing a bond of $\CJ$ but not lying in $\CJ$ has its centre within $\eps/2$ of $\partial\CJ$, where $f=0$. The Neumann condition never enters. (b) $d^*d^*=0$ on plaquette functions, $F(B)=0$. (c) reads off the centres, $f_{\nu\mu}=-f_{\mu\nu}$, and $W^\eps_\mu$ is a segment average of $\sum_\nu\partial_\nu f_{\mu\nu}$. (d) is (c) and the mean value theorem.
\end{proof}
In King's step-$j$ vertex \cite[(3.20)]{K1} this gives $\tau\eta^{3/2}\sum_b5^{-3j}w_\eta(b)(A'_b+A^{(j)}_b)$: one undifferentiated photon end, no explicit power of the spacing, vertex weight $u=\frac52$ in the degree count of \cite[(3.30)]{K1}. By (b) the vertex is invariant under the gauge transformation \cite[(3.46)]{K1} and the constant shift \cite[(3.43)--(3.45)]{K1} drops out of it. Graphs with an odd number of curvature vertices and no external field vanish at $A_j=\phi_j=0$ by the symmetry $(A,\phi)\mapsto(-A,\bar\phi)$.

\paragraph{The degree count on the free $f_D$ flow.} Ba{\l}aban's renormalised degree count \cite[Sect.~3, (3.5)]{Ba3}, in the form King uses on the torus \cite[Sect.~3.5, (3.77)--(3.78)]{K1}: after the rewriting of the divergent subgraphs (Taylor transport of legs, H\"older splits, resummation over orderings up to the lowest external index $j_0$, and evaluation of the resummed coefficients $c_S$ on the free lattice, where they vanish by the counterterm, by sign pairing, by translation invariance or parity, or by the Ward identity of the vacuum polarisation), every ordered ultraviolet stage of a source-free renormalised family has positive degree, and the infrared counterterm lines obey \cite[(3.78)]{K1}. We call this \DEG. It is printed for the $f=1$ flow; the lines of the zone-I graphs are those of the $f_D$ flow, and the transfer is the following.

\begin{lemma}[the two flows differ by a top-scale kernel]\label{lem:a02:RSB}
Let $\DelG_i\defeq G^{f_D}_i-G^1_i$ on $5^{-i}\Z^3$ at $A=0$, in the units of depth $i$. For $a,b\in\{0,1\}$ and every $\gamma<1$, $|\partial^a_x\partial^b_y\DelG_i(x,y)|\le C_\gamma5^{-\gamma i}e^{-\kappa'|x-y|}$, and $\partial_x\DelG_i(\cdot,y)$ is $\alpha'$-H\"older with constant $C_{\alpha'}5^{-s i}$, $s=(1-\alpha')/(1+\alpha')$, times the same decay. In the units of depth $j$ the level-$i$ kernel is the rescaled one, so $\DelG_i$ obeys the slice bounds \cite[(3.63)--(3.65)]{K1} at the index $i$ times $5^{-\gamma i}$: a line with no short-distance singularity.
\end{lemma}
\begin{proof}
Both operators have the same $C^{-1}=-\Delta+m_i^2$, so $\DelG_i=G^{f_D}RG^1$ with $R(u,u')=a_i[1-w(u)w(u')]\mathbf 1[B(u)=B(u')]$. Write $1-w(u)w(u')=(1-w(u))+w(u)(1-w(u'))$; using $\sum_{u\in B(y)}5^{-3i}G^{f_D}(x,u)w(u)=H^{f_D}(x,y)/a_i$ one gets $\DelG_i(x,x')=a_i\sum_yD_0(x,y)g_1(y,x')+\sum_yH^{f_D}(x,y)g_2(y,x')$ with $D_0(x,y)\defeq\sum_{u\in B(y)}5^{-3i}(1-w(u))G^{f_D}(x,u)$, $g_1\defeq\sum_{u'\in B(y)}5^{-3i}G^1(u',x')$, $g_2\defeq\sum_{u'\in B(y)}5^{-3i}(1-w(u'))G^1(u',x')$. \emph{Dipole lemma:} every unit block is a union of finest cells $c$ of side $5^{1-i}$ on each of which $\sum_{u\in c}(1-w(u))=0$ (since $\sum_{u\in c}w=\prod_\mu\sum_{r=0}^4w_1(r)=5^3=|c|$), so $\sum_{u\in B}5^{-3i}(1-w)h(u)=\sum_c\sum_{u\in c}5^{-3i}(1-w(u))[h(u)-h(u_c)]$ is $\le C\sum_c|c|\operatorname{osc}_ch$, and with $|\partial h|\le C(|x-v|+5^{-i})^{-\mathsf p}$ this is $\le C5^{-i}$ for $\mathsf p=2$ and $\le C5^{-i}(1+i)$ for $\mathsf p=3$. The free kernel bounds $|\partial^aG|\le C(|x-u|+5^{-i})^{-1-a}e^{-\kappa|x-u|}$ (\cref{prop:a02:FL}) then give $|D_0|,|g_2|\le C5^{-i}e^{-\kappa\dist}$, $|\partial D_0|,|\partial g_2|\le C5^{-i}(1+i)e^{-\kappa\dist}$, $|g_1|,|\partial g_1|\le Ce^{-\kappa\dist}$, $|\partial^aH^{f_D}|\le Ce^{-\kappa|x-y|}$; every product contains one dipole factor and the $y$-sums converge. For the H\"older clause interpolate the cell oscillation of $h(u)=\partial_xG(x,u)-\partial_xG(\tilde x,u)$ between the H\"older bound $C|x-\tilde x|^{\alpha}r_c^{-2-\alpha}$ and $C5^{1-i}r_c^{-3}$ with $\alpha=(1+\alpha')/2$.
\end{proof}

\begin{theorem}[\DEG\ on the free $f_D$ flow]\label{thm:a02:T}
Every source-free renormalised prefix component $S$ of the free $f_D$ flow, at zero background, admits the rewriting of \cite[Sect.~3]{Ba3} with: (i) every resummed coefficient that vanishes exactly for $f=1$ vanishes exactly for $f_D$; (ii) $c_a^{f_D}=c_a^1+c_a(M_S)$ with $|c_a^1|\le C\rho_{j_0}^{D(S)+|a|}$ the $f=1$ moment and $|c_a(M_S)|\le C_\gamma5^{-\gamma i_0}\rho_{j_0}^{D(S)+|a|}$, $i_0$ the depth of $j_0$; (iii) the same remainder shapes as for $f=1$. Hence \DEG\ holds for the $f_D$ flow, with the degree of every stage unchanged, and the infrared half \cite[(3.78)]{K1} holds with the mass-counterterm blob $U^{f_D}$ differing from $U^1$ by $C(\eps_k/\eta)^\gamma\eta^{1-\varsigma}$ in the units of \cite[(3.78)]{K1}. The $f=1$ input of (ii) on the rows whose coefficient vanishes by the zero-momentum gauge covariance (the Ward rows) is \cref{thm:a02:WR,thm:a02:WRall} of \cref{ssec:a02:ward}; on every other row it is the member-wise identity named there.
\end{theorem}
\begin{proof}
Resum the lines of $S$ to $G^{f_D}_{j_0}$, write $G^{f_D}=G^1+\DelG$ on every line and expand multilinearly: $F_S^{f_D}=F_S^1+M_S$, where $F^1_S$ is the inside family of the same class with $f=1$ lines (the members are the same graphs, since every vertex of a divergent class is weight-free; the counterterms $\dm$, $E_1$ have the same value in both flows; a counterterm piece all of whose lines are $G^1$ is the $f=1$ decomposition of the same counterterm), and $M_S$ is the sum over the members and the non-empty sets $T$ of lines carrying $\DelG$. (i): the exact-zero mechanisms (counterterm per scale assignment, pairing, parity, telescoping of a summable kernel) are identities for every charge-diagonal line kernel when partners carry the same lines, so they hold for the $f_D$ lines directly. (ii): in a term of $M_S$, treat each $\DelG$ line as one line of the fixed index $j_0$; by \cref{lem:a02:RSB} it obeys the index-$j_0$ slice bounds times $5^{-\gamma i_0}$. In any ordering the first $\DelG$ line forces the index $j_0$ on every later line, so the free index sums run over $G^1$ slices only, and the prefixes that contain a $\DelG$ line enter no geometric sum. Each prefix component $\gamma$ with $D(\gamma)\le0$ before the first $\DelG$ line is a source-free, block-free $f=1$ subgraph whose complete inside family (partners included) is present at fixed outside, because replacing $\gamma$ by another member of its class does not touch $T$; renormalise it by the $f=1$ rewriting, finest first, with the $\DelG$ ends as legs at scale $\rho_{j_0}$ (Lipschitz, with $\alpha'$-H\"older first difference, which is what the transport and Taylor steps use), debits landing on a $\DelG$ line paid at $\rho_{j_0}$. The doubly differentiated single line between two current vertices whose other ends carry $\DelG$ is rewritten by the free-lattice identity $\sum D_xD_yG^1(\cdot)=0$ against a constant, whose transport remainder carries $|\DelG(\cdot,v_2)-\DelG(\cdot,v_1)|\le C|v_1-v_2|5^{-\gamma i_0}\rho_{j_0}^{-2}$. After this every prefix with a non-trivial index sum has positive degree, the lines obey their slice bounds, the moment factors $|x_i-x|^{a_i}$ are distributed along paths of lines (costing $\rho_{j_0}^{a_i}$ on a $\DelG$ line), and the multiscale bound (\cref{thm:a02:S3}(a) with the root held) gives $|c_a(M_S)|\le C\rho_{j_0}^{D(S)+|a|}5^{-\gamma i_0|T|}$. (iii) is the same. Infrared half: $U^{f_D}=U^1-(\text{terms with at least one }\DelG_j\text{ soft line})$, and each such term is a skeleton term of the $f=1$ argument in which the soft line obeys the soft-line bounds with the constant $C(\eps_k/\eta)^\gamma$ (\cref{lem:a02:RSB} in physical units), which gives the stated difference.
\end{proof}
% a02_Kpos part deg1: the Ward rows of the degree count on the f = 1 flow (input from a02_Kpos.tex after thm:a02:T)
\subsection{The Ward rows of the degree count}\label{ssec:a02:ward}

The source-free, block-free classes of King's expansion with $D\le0$ are, by order and external profile (\cite[(3.30)--(3.31)]{K1}: $D=\frac o2-3+\frac12(E_S+E_A)+\frac32E_{S_D}$ for a member of order $o$ with $E_S$ undifferentiated scalar, $E_{S_D}$ differentiated scalar and $E_A$ photon legs), the vacuum classes, the photon one- and three-point classes, the scalar two-point classes with and without one photon leg, the photon two-point classes of orders $e^2$, $e^4$, $e^2\lambda$, and the tree \cite[(2.4)]{Ba3}. The resummed coefficient $c_S$ of such a class vanishes, or is coarse-dominated, by one of six mechanisms: the counterterm matched per scale assignment \cite[(3.33)--(3.35)]{K1}; charge parity; a derivative at the end of a line; the locality of the tadpoles; a reflection; and the zero-momentum gauge covariance of the full fixed-order two-photon function and of the photon insertions into the exchange graph $\Sigma_1$. The rows of the last kind, the photon two-point classes (all $24$ configurations, with either photon leg an $A'$-line or an $\tilde A$ field) and the scalar two-point classes with one photon inserted into $\Sigma_1$ at a vertex or on its scalar line, or into the $\lambda$-tadpole loop, are the \emph{Ward rows}. They differ from the others in two respects. Their identity holds for the \emph{whole inside family} with all internal lines resummed and replaced by one kernel, not member by member or at a fixed index assignment; so at a fixed localisation tuple the localisation weights must be common to all members of the family, which under King's rule of one partition of unity per internal vertex \cite[(3.36)]{K1} they are not (at order $e^2$ the bubble has two internal vertices and the $J_2$-loop one, so at a fixed tuple the bubble has no partner, and its weighted zeroth moment grows like $\xi^{-1}$). And in the $R>p$ branch \cite[(3.38)]{K1} King cuts the lines longer than $3\sqrt d$ by vertex positions \cite[p.~666]{K1}, which the identities sum over. Both are met by a reading of King's localisation and cut on data fixed by the tuple alone; this subsection gives it, on the $f=1$ flow, and proves the coarse domination $|c_a|\le C\rho_{j_0}^{D(S)+|a|}$ of \cref{thm:a02:T}(ii) on the Ward rows. Throughout, $S$ is a prefix component of an ordering with $D(S)\le0$, no source, block or $h$-vertex, legs $\Phi_e$ (external fields or lines of index $\ge j_0$ evaluated as fields) of indices $j_e$, $j_0\defeq\min_ej_e$; its \emph{inside family} $F_S$ is the set of all connected block-free subgraphs of the same order and profile occurring in place of $S$ with the same legs at the same outside points, the $\dm$-vertex and $\dm$-inserted graphs included \cite[p.~431]{Ba3}; the \emph{admitted set} $A_S$ is all internal indices $\le j_0$; and for a reference vertex $v_0$ (the vertex of the coarsest leg) and a multi-index $a$ on the other legs the moments are $c_a(x)\defeq\sum_{\text{internal positions}}5^{-3j(V-1)}K_{F_S}(x;x_2,\dots)\prod_{e\ge2}(x_e-x)^{a_e}/a!$, $x$ the position of $v_0$, internal localisations summed.

\begin{lemma}[resummation]\label{lem:a02:RS}
For fixed outside indices, the sum over all orderings of the lines of $S$ and all index assignments in $A_S$ equals the expression with every line of $S$ and of every member of $F_S$ (the internal lines of $\dm$ units included) replaced by the resummed line $G_{j_0}=\sum_{i\le j_0}G_{(i)}$ (resp.\ $G^{\mathrm{ph}}_{j_0}$). Rescaled by $\rho_{j_0}$ to the lattice $\xi\Z^3$, $\xi\defeq5^{-j_0}$, let $C^\xi\defeq(-\Delta+1)^{-1}$ (scalar $\otimes1$; $\delta_{\mu\nu}$ for the photon) and let $F^C_S$ be the family with every line, counterterm-internal lines included, replaced by $C^\xi$ on $\xi\Z^3$. Then for $|a|\le n\defeq-D(S)$
\[
c_a(F_S)=\rho_{j_0}^{D(S)+|a|}\Big[c_a(F^C_S)+\sum_{\emptyset\ne T}c_a(F_{S,T})\Big],\qquad|c_a(F_{S,T})|\le C\ \text{uniformly in }\xi,x,j_0,
\]
where $F_{S,T}$ has the lines of $T$ replaced by $\Delta\defeq G_{j_0}-C^\xi$ (rescaled) or by the periodisation tail on a torus, and the others by $C^\xi$.
\end{lemma}
\begin{proof}
The first statement is \cite[(3.57)--(3.58)]{K1} and multilinearity: King's tie rule is replaced by a simultaneous $S$-first tie-break, which changes nothing in the total. For the second, the resolvent identity $G=C^\xi+G(1-\tilde m^2-\tilde aP)C^\xi$ \cite[p.~437]{Ba3} ($P=Q^*Q$, $\tilde m^2=m^2l_{j_0}^2$, $\tilde a=a_{j_0}$) gives $\Delta=G(1-\tilde m^2-\tilde aP)C^\xi$; with $|G(x,y)|\le C|x-y|_\xi^{-1}e^{-\delta|x-y|}$, $|DG|\le C|x-y|_\xi^{-2}e^{-\delta|x-y|}$, $|D_xD_yG|\le C|x-y|_\xi^{-3}e^{-\delta|x-y|}$ (\cref{prop:a02:FL}; \cite[Thm.~3.3]{K1}), the same for $C^\xi$, and $P$ bounded with unit-range kernel, convolution in $d=3$ gives $|\Delta|\le C$, $|D\Delta|\le C(1+\ln_+(1/|x-y|_\xi))$, $|D_xD_y\Delta|\le C|x-y|_\xi^{-1}$, all times $e^{-\delta'|x-y|}$: a $\Delta$-line is a unit-scale kernel, undifferentiated bounded and Lipschitz below scale $1$, differentiated two orders better than the $G$-line (degree $+2$ up to a logarithm absorbed by an $\epsilon$-loss). The replacement is multilinear, the $\dm$ partners' internal lines being expanded by the same rule, so every all-$C$ subgraph of $F_{S,T}$ is renormalised inside $F_{S,T}$ by its own partners. On a torus $C^\xi_T=\sum_\gamma C^\xi(\cdot+\gamma)$ and the $\gamma\ne0$ tail is smooth and $O(e^{-\delta|\gamma|})$, one more $\Delta$-type line. Decompose the $C^\xi$ lines into scales ($C^\xi$ is the $A=0$, mass-$1$ instance of \cite[(2.13)--(2.17)]{K1} on $\xi\Z^3$ with the bounds \cite[(3.63)--(3.65)]{K1}) and order all lines: a prefix component of an output term is either all-$C$, hence a proper subgraph of a member of $F_S$ (lower order, or the same order and more legs), whose class is rewritten by the same statement for the $C^\xi$ line system (induction on order, then on the number of legs, which is well founded since $E\le4$ for every divergent class); or it contains a $\Delta$-line: undifferentiated, it sits at the top scale and enters no fine-scale sum; differentiated, its degree is $\ge D_{\mathrm{plain}}+2-\epsilon\ge1-\epsilon>0$, because every non-vanishing non-vacuum source-free connected subgraph has $D\ge-1$ (the $D=-2$ classes are the current tadpole, which vanishes by charge parity, and vacuum graphs, which are not prefix components of a non-vacuum $S$). So every ordered subgraph of every output term has positive degree, the multiscale bound \cref{thm:a02:S3}(a) applies, and the final scale sum is coarse-dominated.
\end{proof}

\begin{lemma}[zero-momentum gauge covariance]\label{lem:a02:GC}
Let the lines be gauge covariant in a background $B$, $C_B(x,y)\mapsto g(x)C_B(x,y)g(y)^T$ under $B\mapsto B+\nabla\chi$, $g=e^{eq\chi}$, with King's vertices \cite[(3.18)--(3.23)]{K1} at background $B$, and on $\xi\Z^3$ let $B$ be constant, $\chi_B(x)=B\cdot x$. For a labelled contraction $G$ with two external scalar fields at $x,x_2$, $K^G_B(x,x_2)=g(x)K^G_0(x,x_2)g(x_2)^T$, and for $G$ with photon legs only $F^G_B=F^G_0$. Differentiating at $B=0$ in direction $\mu$: \textup{(a)} with $W(A)\defeq\ln Z(A)-\ln Z(0)$, $Z(A)=\int d\mu(A_{\mathrm{int}})\int d\mu_{C_0}(\phi)e^{V(\phi,A+A_{\mathrm{int}})}$, $W(A+B)=W(A)$ for constant $B$ on $\xi\Z^3$, hence for the full fixed-order connected two-photon function $\Pi=\partial^2W/\partial A\partial A|_0$, all its labelled contractions and $\dm$-inserted ones included, $\sum_{b_2\in\mathrm{bonds}_\nu}\Pi_{\mu\nu}(b_1,b_2)=\partial_{A_{b_1}}\partial_{B_\nu}W(A+B)|_0=0$; \textup{(b)} for every two-scalar-leg parent $G$ with scalar charge structure, $\sum_{b\in\mathrm{bonds}_\mu}K^{(G+\mu)}(x,x_2;b)=eq(x-x_2)_\mu K^G_0(x,x_2)$, so the zeroth moment of the insertion family of $G$ equals $-eq$ times the first moment of $G$. The insertions are: on a scalar line $-C(\partial_B(-\Delta_B))C$, the vertex $J_1$; at a vertex $J_m(B)$ through $U(B_b)$ a $J_{m+1}$-type vertex and through $D_B$ a $Q_{m+1}$-type vertex; none on photon lines, $\lambda$ or $M$.
\end{lemma}
\begin{proof}
(a) $Z(A+\nabla\chi)=Z(A)$ by the change of variables $\phi\mapsto e^{eq\chi}\phi$, a constant $B$ is $\nabla(B\cdot x)$ on $\xi\Z^3$, and the order-by-order expansion of $W$ in $e$ commutes with the shift because $W$ depends on $eA$ only. (b) Per contraction the value is a contraction of covariant tensors with covariant propagators, the internal phases cancel, $\partial_{B_\mu}g=eqx_\mu g$, and $K^G_0$ commutes with $q$; the insertion list is Leibniz on $C_B$ and on $J_m(B)$ \cite[(3.22)]{K1} ($Q_m$ carries no $U$ and is $B$-independent).
\end{proof}
The identity is used only for $F^C_S$ of \cref{lem:a02:RS}, whose lines are full resummed lines replaced by $C^\xi$ (the total-field form); no identity is applied to a fixed ordering or a fixed index assignment. On a torus (a), (b) hold up to the holonomy response $O(e^{-\delta|T|_\xi})$, one more $\Delta$-type term.

\begin{lemma}[reflection]\label{lem:a02:REF}
Let $g_s$, $c$ be the even scalar and photon kernels $C^\xi$ on $\xi\Z^3$. \textup{(a)} The exchange kernel $\Sigma_1$ with both $S_D$ legs on the scalar line and $\mu=\nu$ is $k(y)=-e^2\sum_\mu(\Delta_\mu g_s)(y)c(y)$, $y=x-x_2$, which is even, so its first moment vanishes. \textup{(b)} For the order-$e^2$ two-photon family ($J_1J_1$ bubble, $J_2$ loop, $Q_2$ loop) with zeroth moment $M^{(0)}=0$, every first moment $M^{(1)}_{\mu\nu\sigma}\defeq\sum_{b_2}\Pi_{\mu\nu}(b_1,b_2)(b_2-b_1)_\sigma$ vanishes.
\end{lemma}
\begin{proof}
(a) With $b=(x-\xi e_\mu,x)$, $b'=(x_2-\xi e_\mu,x_2)$, $\partial_b\partial_{b'}g_s(x-x_2)=-\Delta_\mu g_s(x-x_2)$ and the photon line is $c(x-x_2)$. (b) Under the reflection $R_\sigma$ the bond $(z,\mu)$ goes to $(Rz-\xi e_\sigma\delta_{\mu\sigma},\mu)$ and $A_\mu\mapsto s_\mu A_\mu$ ($s_\sigma=-1$, else $+1$); the family is the expansion of a reflection-invariant action, so $\Pi_{\mu\nu}(b_1,b_2)=s_\mu s_\nu\Pi_{\mu\nu}(R_\mu b_1,R_\nu b_2)$, hence $M^{(1)}_{\mu\nu\sigma}=s_\mu s_\nu[-M^{(1)}_{\mu\nu\sigma}-\xi(\delta_{\nu\sigma}-\delta_{\mu\sigma})M^{(0)}_{\mu\nu}]$ under $R_\sigma$ and $=s'_\mu s'_\nu M^{(1)}_{\mu\nu\sigma}$ under $R_\tau$, $\tau\ne\sigma$. With $M^{(0)}=0$: if $s_\mu s_\nu=+1$ under $R_\sigma$ then $M^{(1)}=0$; otherwise exactly one of $\mu,\nu$ is $\sigma$, say $\mu=\sigma\ne\nu$, and $R_\nu$ gives $M^{(1)}=-M^{(1)}$.
\end{proof}

\paragraph{The localisation by centres.} King localises every graph of $P^{(j),1}$ by inserting $1=\sum_z\chi_z$ at each internal vertex \cite[(3.36)]{K1} and summing over the tuple $(\{y_i\},\{z\})$ \cite[(3.37)]{K1}. We attach the partitions of unity to \emph{centres} instead: \textup{(i)} one centre to every $A'$ field of a $J_m$ or $Q_m$ vertex ($m$ centres at $J_m$ \cite[(3.22)]{K1}, $m$ at $Q_m$ \cite[(3.23)]{K1}); \textup{(ii)} one centre to every other vertex ($\lambda$ \cite[(3.18)]{K1}, block vertices \cite[(3.24)--(3.26)]{K1}, $h$-, $g$- and source vertices). For each centre $c$ insert $1=\sum_{z_c}\chi_{z_c}(p_c)$, $p_c$ the position of the vertex carrying $c$; the tuple of a graph is $\tau=(\{y_i\},(z_c)_c)$; $\square$ is the smallest cube of large blocks containing every $y_i$ and every $z_c$ (the centres of the localised counterterm graphs included), and $\square'$, $\square''$, $x_0$, $B=A_j(x_0)$, $\tilde A$ and $R\defeq\dist(\{y_i\},\{z_c\})$ are as in \cite[pp.~660--662]{K1}. King's index rules \cite[p.~667]{K1} are kept: a $\dm$ graph is localised only if all its integers are $<j$, and an $E_1$ piece with $m_1<m$ gets one centre at one vertex connected to $l(m)$. King's printed rule is this rule with one centre per vertex in place of (i).

\begin{lemma}[King's localisation argument with centres]\label{lem:a02:Lsharp}
\textup{(a)} The only properties of \cite[(3.36)]{K1} used in \cite[Sects.~3.4--3.5]{K1} are: $\sum_z\chi_z=1$ per inserted factor; $0\le\chi\le1$ with support of size $O(1)$ and bounded derivatives; every localised vertex lies within $O(1)$ of a centre in $\square$; and the number of centres is bounded by a function of $\nbar$ and the order. The rule by centres has all of them, so \cite[Thm.~3.5, Prop.~3.6, (3.40)--(3.42), (3.43)--(3.57), (3.68), (3.85)]{K1} and the step comparison of \cite[p.~665]{K1} hold verbatim, with constants depending on $\nbar$ and $d$ (the centres of one vertex lie within $\frac43\sqrt d$ of each other, so the tree distance over centres exceeds that over vertices by at most $C(\nbar)$). \textup{(b)} \emph{Centre count.} For $S'\in F_S$ with $e$-order $\alpha$, $\lambda$-order $\beta$ and $E_{\tilde A}$ legs that are $\tilde A$ fields, the number of centres is $N(S')=\alpha-E_{\tilde A}+\beta$; it is the same for every member of $F_S$. \textup{(c)} Consequently every graph of a renormalisation orbit carries the same tuple, under a fixed bijection of the centres of $S$ and $S'$, hence the same $\square$, $\square'$, $x_0$, $B$, $\tilde A$ and the same branch of \cite[(3.38)/(3.39)]{K1}; the same along any sequence of replacements, nested units included. \textup{(d)} In the order-$e^2$ photon two-point family every centre sits at the $S$-end of an $A'$-line leg: no member has an internal photon line or a $\lambda$ vertex, and two $\tilde A$ legs are impossible.
\end{lemma}
\begin{proof}
(a) Each inserted factor sums to $1$, so \cite[(3.37)]{K1} is an identity for any number of centres; the uses on \cite[pp.~659--668]{K1} are: the definition of the tree distance and the decay $\exp[-\delta\dist]$ (co-located centres cost a bounded factor); the dichotomy and $\sum_{y,z}e^{-\frac12\delta R}\le C\eta^{-3}|\CJ|$; the cubes; $B=A_j(x_0)$ on $\square'$; the gauge transformation and $\theta$; $|x_0-z_c|\le Cp$ by the tree decay; \cite[(3.68)]{K1}, the weight of a vertex being a product of $\chi$'s in $[0,1]$ supported in a unit cube; the root sum by $|\square'|$; the replacement of $\chi(x')$ by $\chi(x)$ in the step comparison, since products of $\chi$'s have bounded derivatives; $\delta G$ bounds ``since $x,y\in\square$'', every vertex lying within $\frac23\sqrt d$ of each of its centres; the counterterm localisation with the index rules; and the fold \cite[(3.85)]{K1} of a leg's $\chi$ into the leg. (b) Each $A'$ field at a $J$- or $Q$-vertex carries one factor $e_j$ \cite[(3.22)--(3.23)]{K1}; each $\tilde A$ carries one $e_j$ \cite[(3.47)--(3.53)]{K1} and is never contracted; each $\lambda$ vertex carries one $\lambda_j$. So $\alpha=\sum_v(a_v+t_v)$ with $a_v$ the $A'$ fields and $t_v$ the $\tilde A$ legs at $v$, $E_{\tilde A}=\sum_vt_v$, $\beta=|\Lambda(S')|$, and $N=\sum_va_v+|\Lambda(S')|$. $F_S$ fixes $(\alpha,\beta)$ and the legs; every identity used acts within one $(\alpha,\beta)$, since gauge invariance holds for every value of $e$ and $\lambda$ and $\dm$, $E_1$ are sums of graphs of definite order \cite[(3.33)--(3.35)]{K1}. (c) The reduction is graph-local and the listed objects are functions of the tuple. (d) $\alpha=2$, $\beta=0$, $E_{A'}+E_{\tilde A}=2$, so there are exactly $E_{A'}$ $A'$ fields, the $S$-ends of the $A'$-line legs; a $J$/$Q$ vertex needs an $A'$ field, and a subgraph of $\tilde A$-only vertices is the tree \cite[(2.4)]{Ba3}, not a two-photon graph.
\end{proof}

\begin{lemma}[frozen weights]\label{lem:a02:FW}
Let $S$ be a Ward-row unit at a fixed tuple, $x$ its reference point, and for a member $S'$ let $W_{S'}=\prod_{c\in\mathrm{legC}}\chi_{z_c}(p_c)\prod_{c\in\mathrm{intC}}\chi_{z_c}(p_c)$, where $\mathrm{legC}$ are the centres at the $S$-ends of the $A'$-line legs (common to all members) and $\mathrm{intC}$ the remaining $N-E_{A'}$ centres (common by \cref{lem:a02:Lsharp}(b), identified by the fixed bijection). Put $\tilde\Phi_\ell\defeq\chi_{z_\ell}\Phi_\ell$ for each $A'$-line leg, $\tilde\Phi\defeq\Phi$ for the other legs, and $W^\circ\defeq\prod_{c\in\mathrm{intC}}\chi_{z_c}(x)$, a member-independent number. Then
\[
\sum_{S'\in F_S}\sum_{\mathrm{pos}}K_{S'}W_{S'}\prod_\ell\Phi_\ell=W^\circ\sum_{S'}\sum_{\mathrm{pos}}K_{S'}\prod_\ell\tilde\Phi_\ell+\sum_{S'}\sum_{\mathrm{pos}}K_{S'}\Big[\prod_{c\in\mathrm{intC}}\chi_{z_c}(p_c)-W^\circ\Big]\prod_\ell\tilde\Phi_\ell.
\]
The first term is the unweighted family sum against the legs $\tilde\Phi$ (unit-scale smooth factors times legs, H\"older or $C^{1,\alpha}$ at the leg's own scale, \cite[(3.85), p.~665]{K1}), times a constant. The second term is a sum of remainder terms with a credit $\rho_{j_l}$ on a line $l$ of $S$ and a debit on a unit-scale weight (cost $\rho_j^{-1}=1$), so $D'(S)=D(S)+1$ and $D'(\gamma)\ge D(\gamma)$ for $\gamma\subsetneq S$; it is absent for the order-$e^2$ photon two-point class ($\mathrm{intC}=\emptyset$ by \cref{lem:a02:Lsharp}(d)), the only Ward row with $D=-1$, where the $n=1$ Taylor step acts on $\tilde\Phi$, King's fold \cite[(3.85)]{K1}.
\end{lemma}
\begin{proof}
Add and subtract; for products of $[0,1]$-valued functions $|\prod a_i-\prod b_i|\le\sum|a_i-b_i|$, so $|\prod\chi(p_c)-W^\circ|\le\sup|\nabla\chi|\sum_{c}|p_c-x|\le C\sum_{l\in P}|\mathrm{step}_l|$ along a path $P$ of lines of $S'$ from $x$ to the vertices carrying $\mathrm{intC}$; each step is absorbed by part of its line's decay at the price $\rho_{j_l}$ \cite[(3.13)--(3.14), (3.19)]{Ba3}. Every Ward row other than the order-$e^2$ photon class has $D(S)=0$, hence $D'(S)=1>0$.
\end{proof}

% a02_Kpos part deg2: the Ward rows in the R > p branch (the cut by clusters), and the conclusion for (DEG) (input after a02_deg1)
\begin{theorem}[the Ward rows, $R\le p$]\label{thm:a02:WR}
Let the localisation be by centres, and let $S$ be a Ward-row prefix component at a fixed tuple in the branch $R\le p$ of \cite[(3.39)]{K1}, after King's reduction \cite[(3.43)--(3.55)]{K1} with the gauge $B=A_j(x_0)$ of the cube. Then the sum over $F_S$ and $A_S$ is a sum of terms of three shapes: \textup{(R0)} zero; \textup{(R1)} local terms $\sum_x5^{-3j}c_a(x)\prod\partial^a\tilde\Phi(x)$, $|a|\le n=-D(S)$, with $|c_a(x)|\le C\rho_{j_0}^{D(S)+|a|}$ uniformly in $x$, $j_0$, $k$, $\eta$ and the tuple, times $W^\circ$ and the frozen $\tilde A(x)$-legs; \textup{(R2)} remainders with a credit $\rho^\theta$ on a line of $S$ and a debit of order $\le\theta$ on one leg of $S$ or on a unit-scale weight, $D'(S)>0$ and $D'(\gamma)\ge D(\gamma)$ for every prefix component $\gamma\subsetneq S$. The constant depends on $\nbar$, $d$, $L$, $\alpha$ only.
\end{theorem}
\begin{proof}
(1) By \cref{lem:a02:Lsharp}(c) every member of $F_S$, and every member of every family met in the nesting, is reduced with the same $\square$, $\square'$, $x_0$, $B$, $\tilde A$ and lies in the same case below. (2) Insert at every $\tilde A$-vertex $v$ the identity $1=\zeta(v)+(1-\zeta(v))$, $\zeta(v)\defeq\psi(\dist(v,\square)/p)$, $\psi\in C^\infty$, $\psi=1$ on $[0,1]$, $0$ on $[\frac32,\infty)$, $|\psi'|\le4$. An $\tilde A$-vertex, an insertion vertex of the expansion \cite[(3.52)--(3.53)]{K1} of one propagator in $\tilde A$, has exactly two scalar line-ends, both on the chain of that propagator, and external $\tilde A$-legs only; so a connected subgraph all of whose vertices are $\tilde A$-vertices is a path of $r$ vertices inside one chain, with $E_S+E_{S_D}=2$, $E_A\ge r$ and order $\ge r$, hence $D\ge\frac r2-3+\frac12(2-E_{S_D}+r)+\frac32E_{S_D}=r-2+E_{S_D}$ by the degree formula of \cref{ssec:a02:ward}, and $D\le0$ only for $r=2$, $E_{S_D}=0$, $E_A=2$: the tree \cite[(2.4)]{Ba3}, which is not a Ward row (its member-wise identity at every tuple, the cut of its position sums at $\partial\square'$ included, is \cref{prop:a02:TREE} below). So $S$ contains a vertex $u$ of the original graph, localised by its centre within $\sqrt d$ of $\square$. A term in which some $\tilde A$-vertex $v$ carries $1-\zeta$ has $\dist(v,\square)\ge p$, hence $|u-v|\ge p-\sqrt d$; a term with $\dist(x,\square)>p/2$ has $|x-u|\ge p/2-\sqrt d$; a term with a line $\delta G_{(i)}(\square')$ has either both ends of that line at distance $\ge p/2$ from $\partial\square'$, or an end $w$ with $\dist(w,\square)\ge\frac32p$, hence $|w-u|\ge\frac32p-\sqrt d$. Two vertices of $S$ at distance $\ge t$ are joined by a path of at most $V-1$ lines of $S$, $V\le C(\nbar)$ its number of vertices, so some line $\ell$ of $S$, of index $i$, has length $\ge t/V$, and $e^{-\delta_0|\ell|/\rho_i}\le C_\theta(V\rho_i/(\delta_0t))^\theta e^{-\delta_0|\ell|/(2\rho_i)}$ for every $\theta\ge0$ (the configuration space is split by the first line of the path of length $\ge t/V$; the split is sharp but harmless, since no identity is used on these terms); a $\delta G_{(i)}(\square')$ line with both ends at distance $\ge p/2$ from $\partial\square'$ carries $e^{-\delta_1p/(2\rho_i)}$ by \cref{lem:a02:I}. In every case the line supplies a credit of any order, $\le C_\theta(\rho_i/p)^\theta$ with $\theta=3>-D(S)$, on a line of $S$ with no debit and no identity (\cref{prop:a02:FL}, \cref{lem:a02:I}); these terms are of shape (R2), bounded with absolute values, and $S$ needs no rewriting. (3) In the remaining \emph{bulk} case ($\dist(x,\square)\le p/2$, every line a free slice, every $\tilde A$-vertex carrying $\zeta$, so $\zeta(x)=1$ and all position sums are sums over the whole lattice) \cref{lem:a02:FW} turns the weighted family into $W^\circ$ times the unweighted family against $\tilde\Phi$ plus (R2) weight units, and the unweighted family is the zero-background free-lattice object of \cref{lem:a02:RS}: resum and replace by $C^\xi$; for the photon two-point classes $F^C_S$ is the full fixed-order connected two-photon function of the $C^\xi$ theory, all ways of filling the inside and the $\dm$ units included \cite[p.~431]{Ba3}, so $c_0(F^C_S)=0$ by \cref{lem:a02:GC}(a), and at order $e^2$ also $c_1(F^C_S)=0$ by \cref{lem:a02:REF}(b); for the insertion classes the inside families are exactly the insertion families of the order-$2$ parents $\Sigma_1$, the $\lambda$-tadpole and the $Q_2$-tadpole, so $c_0(F^C_S)=-eq\times(\text{first moment of the parent})^C=0$ by \cref{lem:a02:GC}(b), \cref{lem:a02:REF}(a) and the locality of the tadpoles; \cref{lem:a02:RS} bounds the remaining terms by $C\rho_{j_0}^{D+|a|}$. The transport and Taylor steps \cite[(3.10), (3.34)]{Ba3}, with the reference vertex the vertex of the coarsest leg and the Taylor leg the finer one, then give (R1) and (R2) with the degree arithmetic of \cite[Sect.~3]{Ba3}: each remainder factor $|x_e-x|^\theta$ is bounded along a path of lines of $S$ at the price $\rho_{j_l}^\theta$, the H\"older quotient of the leg costs $\rho_{j_e}^{-\theta}$ \cite[(3.65), (3.85)]{K1}, $D'(S)=D(S)+\theta>0$, and a proper prefix component does not contain the debited leg. (4) Sub-units are rewritten first, finest first, at the same tuple; each rewriting is an identity on its own family, so the moments of $S$ are those of its original weighted family kernel, to which (3) applies; a sub-unit's frozen weights are frozen again at $x$, since products of $[0,1]$-valued unit-scale factors compose; each leg receives at most one unit's debit, of order $\le1+\alpha$; a weight debit is never a leg debit; and one gauge, $B=A_j(x_0)$ per cube, is used throughout.
\end{proof}
% a02_Kpos part deg2, continued: the tree near the face of the reduction cube (input from a02_deg2.tex after thm:a02:WR)
\paragraph{The tree near the face of the reduction cube.} Step (2) of the proof of \cref{thm:a02:WR} shows that the only source-free connected subgraph with $D\le0$ all of whose vertices are $\tilde A$-vertices is the tree $T$ of \cite[(2.4)]{Ba3}: two insertion vertices of \cite[(3.52)--(3.53)]{K1} of $J_1$ type joined by a doubly differentiated line, $D(T)=0$, no partner, inside family $\{T\}$. Its member-wise identity, a derivative at the end of a line (\cref{ssec:a02:ward}), is a sum over the whole lattice; at a fixed tuple the positions of insertion vertices are summed over the reduction cube $\square'$ with a sharp cut (\cite[p.~668]{K1}: only through the expansion \cite[(3.52)]{K1} can vertices approach $\partial\square'$), so near $\partial\square'$ the identity leaves a face term. The lines are decomposed scale by scale on $\square'$, $G_{(i)}(\square',0)=G_{(i)}(0)+\delta G_{(i)}(\square')$, the difference kept as an internal line of index $i$ (as in step (2) of the proof of \cref{thm:a02:WR}); we write $F_i(v)\defeq\exp[-\delta_1\dist(v,\partial\square')/\rho_i]$ with $\delta_1\defeq\delta_0/2$, $\delta_0$ the decay rate of the free slices \cite[(3.63)]{K1}.

\begin{lemma}[face factors]\label{lem:a02:FACE}
\textup{(a)} For $u,v\in\square'$ the line $\delta G_{(i)}(\square')(u,v)$ obeys the free slice bounds \cite[(3.63)--(3.65)]{K1} (sup, first differences, H\"older quotients of the first differences) at index $i$ with half the rate, times $F_i(u)F_i(v)$ (for a H\"older pair $u,\tilde u$, $F_i$ taken at the point of the pair nearer to $\partial\square'$). \textup{(b)} Let a term of the ordered bound of \cref{thm:a02:S3}(a) carry the factor $F_i(v)$ at a vertex $v$ of the tree $T$, $i$ the index of its line $\ell$. Then the term is an \textup{(R2)} unit in the sense of \cref{thm:a02:WR}: a credit $\rho_i$ on $\ell$ and a debit $\rho_{i'}^{-1}$ on one leg $l'$ of $T$, of index $i'\ge i$, or no debit when the surviving vertex of $T$ is the root of the graph; $D'(T)=1$, and $D'(\gamma)\ge D(\gamma)$ for every prefix component $\gamma\supsetneq T$.
\end{lemma}
\begin{proof}
(a) The boundary clause of \cite[Thm.~3.3]{K1} gives $\delta G_j(\square')$ and $\delta C^{(j)}(\square')$ their bounds with the extra factor $e^{-\delta_0\dist(\{u,v\},\partial\square')/\rho_j}$ (the clause used for the lines $\delta G_j(\square',A)$ in the proof of \cref{lem:a02:I} through \cite[(5.3)]{K2}). Writing each factor of the slice $G_{(i)}=a_iG_iQ^*C^{(i)}QG_i$ on $\square'$ as its free-lattice part plus its $\delta$-part and telescoping, every term of $\delta G_{(i)}(\square')$ carries one such factor, and the decay $e^{-\delta_0|\cdot|/\rho_i}$ of the factors (the $\delta$-factor's own included) moves it to both $u$ and $v$ at half the rate ($\min+\text{distance}\ge\max\ge\frac12(\dist(u,\partial\square')+\dist(v,\partial\square'))$), which is $F_i(u)F_i(v)$ with $\delta_1=\delta_0/2$ against the free bounds at the rate $\delta_0/2$; the H\"older clause follows from the sup and first-difference clauses of the factors as \cite[Prop.~3.7]{K1} follows from \cite[Thm.~3.3]{K1}, the boundary factor of the differentiated factor taken at the point of the pair nearer to $\partial\square'$, a pair at distance $>\rho_i$ being bounded by two first differences. On the $f_D$ flow the same follows from the image sum $\delta G_{(i)}(\square')(u,v)=\sum_{g\ne1}\operatorname{sgn}(g)\,G_{(i)}(0)(u,gv)$ of \cref{prop:a02:S}(c), since $|u-gv|\ge\max\{|u-v|,\dist(u,\partial\square')+\dist(v,\partial\square')\}$ for every $g\ne1$ and the images at distance $\le t$ number $O((1+t/p)^d)$. (b) In the ordered bound the line $\ell=(v,v')$ is shrunk by summing one endpoint against its kernel; $\dist(\cdot,\partial\square')$ is $1$-Lipschitz, so $F_i(v)\le e^{\delta_1|v-v'|/\rho_i}F_i(v')$ moves the factor to whichever endpoint survives, at the cost of half of the decay of $\ell$. The surviving vertex $x$ of $T$ is then summed against a leg $l'$ of $T$, a line of index $i'\ge i$ with the kernel profile $|K_{l'}(u-x)|\le C\rho_{i'}^{-c}e^{-\delta_0|u-x|/\rho_{i'}}$ of \cite[(3.63)--(3.65)]{K1} (a cut-line leg, \cref{lem:a02:XL}(b), has the unit-scale profile, $i'=j$), or it is the root. Since $F_i(x)\le\sum_\pi e^{-\delta_1\dist(x,\pi)/\rho_i}$ over the $2d$ face planes of $\square'$, each term depending on one coordinate only, summing that coordinate first gives $\sum_{x_\mu}5^{-j}e^{-\delta_1|x_\mu-\pi_\mu|/\rho_i}\le C\rho_i$ (as $5^{-j}\le\rho_i$), while the sum of $|K_{l'}|$ over the other $d-1$ coordinates is at most $C\rho_{i'}^{-c+d-1}$; so
\[
\sum_{x}5^{-3j}|K_{l'}(u-x)|\,F_i(x)\ \le\ C\,\frac{\rho_i}{\rho_{i'}}\,\rho_{i'}^{\,d-c},
\]
the bound $C\rho_{i'}^{d-c}$ of that step of the induction times $\rho_i/\rho_{i'}$. At the root, where the final position sum is paid by $|\square'|$ (\cite[p.~664]{K1}), $\sum_{x\in\square'}5^{-3j}F_i(x)\le C\rho_i\cdot(\text{area of }\partial\square')\le C(\rho_i/p)\,|\square'|$, the side of $\square'$ being $\ge4p$. In the telescoped product of the ordered bound (\cref{lem:a02:Dsp}(ii)) the factor $\rho_i/\rho_{i'}$ raises by one the exponent of exactly those prefixes that contain $\ell$ and not $l'$, so $D'(T)=D(T)+1=1$ and $D'(\gamma)\in\{D(\gamma),D(\gamma)+1\}$ for $\gamma\supsetneq T$; at the root the credit has no debit.
\end{proof}

\begin{proposition}[the tree at every tuple]\label{prop:a02:TREE}
Let $T$ be the tree \cite[(2.4)]{Ba3} as a prefix component at a fixed tuple, in either branch (with $\square'_K$ in place of $\square'$ in the branch $R>p$), with vertices $v_1$, the reference point $x$, and $v_2$, both summed over $\square'$, its line $\ell$ of index $i$, and the products $W(v_1)$, $W'(v_2)$ of the factors at its vertices (an $\tilde A$ field, the weight $\zeta$ or $1-\zeta$ of step (2) of the proof of \cref{thm:a02:WR}, and one chain-end line of index $\ge i$, a cut line in the branch $R>p$). Then $\sum_{v_1,v_2\in\square'}\ell(v_1,v_2)W(v_1)W'(v_2)$ is a sum of terms of the shapes \textup{(R0)}--\textup{(R2)} of \cref{thm:a02:WR}, with $D'(T)\ge\min\{\alpha,1\}>0$ and $D'(\gamma)\ge D(\gamma)$ for $\gamma\supsetneq T$. In particular the member-wise identity of the tree holds at every tuple, near $\partial\square'$ included, which is what \cref{cor:a02:DEGall}(a) uses for it.
\end{proposition}
\begin{proof}
If $\ell=D_{b_1}D_{b_2}G_{(i)}(0)$, write
\[
\sum_{v_1,v_2\in\square'}\ell\,W\,W'=\sum_{v_1\in\square'}W(v_1)\sum_{v_2\in\square'}\ell(v_1,v_2)\big[W'(v_2)-W'(v_1)\big]+\sum_{v_1\in\square'}W(v_1)W'(v_1)\Big[\sum_{v_2}\ell(v_1,v_2)-\sum_{v_2\notin\square'}\ell(v_1,v_2)\Big].
\]
The full sum over $v_2$ vanishes exactly, a lattice difference summed over the whole lattice (or the torus) being zero, the derivative at the end of a line: \textup{(R0)}. The first term is the transport step of \cite[(3.10)]{Ba3} with $n=0$: the factor $|v_1-v_2|^\alpha$ is absorbed by the decay of $\ell$ at the price $\rho_i^\alpha$, a credit on $\ell$, and the $\alpha$-H\"older quotient of $W'$ costs $\rho_{i'}^{-\alpha}$ on the chain-end line at $v_2$, which is H\"older at its own scale up to $\partial\square'$ (the free part by \cite[(3.65)]{K1}, the $\delta G$ part by \cref{lem:a02:FACE}(a), a cut line by \cref{lem:a02:XL}(b)), and $O(1)$ on the smooth factors $\tilde A$ (\cite[(3.45)]{K1}) and $\zeta$ ($|\nabla\zeta|\le4/p$): \textup{(R2)} with $D'(T)=\alpha$. In the cut term $v_1\in\square'$ and $v_2\notin\square'$, so $|v_1-v_2|\ge\dist(v_1,\partial\square')$ and half of the decay of $\ell$ gives $F_i(v_1)$; the sum over $v_2$ is then bounded by the kernel integral of $\ell$, as at every step of the ordered bound, and \cref{lem:a02:FACE}(b) makes the term an \textup{(R2)} unit with credit $1$. If $\ell=\delta G_{(i)}(\square')$, \cref{lem:a02:FACE}(a) gives $F_i(v_1)$ (the factor $F_i(v_2)$ is dropped) and \cref{lem:a02:FACE}(b) applies; no identity is used. Since $T$ has no partner and its inside family is $\{T\}$, the rewriting is an identity term by term, and the nesting of step (4) of the proof of \cref{thm:a02:WR} is unchanged. Away from $\partial\square'$ and on the torus this is Ba{\l}aban's treatment of \cite[(2.4)]{Ba3} in \cite[Prop.~2.1]{Ba3}; the cut term is the face term that the cut of the position sums produces.
\end{proof}


\paragraph{The cut by clusters.} In the branch $R>p$ King proves \cite[(3.38)]{K1} by treating every propagator longer than $3\sqrt d$ as an external line, so that the graph factorises into localised subgraphs, each reduced with its own cube \cite[p.~666]{K1}. That cut is by vertex positions: at a fixed tuple the partition into subgraphs, hence each subgraph's $\square$, $x_0$, $B$ and $\tilde A$, varies with the positions that the identities of \cref{ssec:a02:ward} sum over, within one member as well as across members. We cut by the centres instead. Units of depth $j$, $d=3$; $c_0\defeq\frac{13}3\sqrt d$ and $R_0\defeq c_0-\frac43\sqrt d=3\sqrt d$, King's threshold. Fix a graph $H$ of $P^{(j),1}$ localised by centres and a tuple $\tau$ with point set $P(\tau)$ (the $y_i$ and the $z_c$).
\begin{enumerate}[label=\textup{(C\arabic*)},nosep,leftmargin=*]
\item \emph{Clusters.} Let $\Gamma_\tau$ be the graph on $P(\tau)$ with an edge between two points at Euclidean distance $\le c_0$; the clusters $K$ are the components of $\Gamma_\tau$, with point sets $P_K$.
\item \emph{Vertices.} A vertex with centres is assigned to the cluster of its centres (well defined on the support of the weights, (a$'$) below); a block or $y$-point carries its own point of $P(\tau)$. An unlocalised vertex is either the attachment vertex of a $\dm$ point vertex whose blob carries an index $\ge j$ (\cite[p.~667]{K1}: no partition of unity is introduced in that graph) or the vertex of a one-vertex graph; the former has exactly two scalar legs \cite[(3.19)]{K1}, and each maximal path through such vertices is read as one \emph{dressed line} between localised ends; one-vertex graphs have no internal line, no divergent subgraph and no cube, and are bounded by the decay of their external lines; an $E_1$ piece localised at one vertex forms one cluster.
\item \emph{Cut.} A line (scalar, photon, external line $a_jG_jQ_j^*$ to a $y$-point, fine $\dm$-blob line or dressed line) is \emph{cut} iff its end-vertices lie in different clusters. A cut line is kept whole, with no scale decomposition, and treated as an external line of both clusters; the \emph{cluster graph} $H_K$ is $H$ restricted to the vertices of $K$, with its uncut lines internal and its cut lines as external fields of index $j$.
\item \emph{Reduction and families.} Each $H_K$ is reduced by \cite[(3.43)--(3.55)]{K1} with its own cube: $\square_K$ the smallest cube of large blocks containing $P_K$, $\square'_K$, $\square''_K$ its $2p$- and $p$-neighbourhoods, $x_0^K$ the centre of $\square_K$, $B_K=A_j(x_0^K)$, $\tilde A_K=A^{(j)}-B_K$ on $\square'_K$; cut lines keep the background $A^{(j)}$ and receive King's unitary factors \cite[(3.46)]{K1} at their ends, as \cite[p.~661]{K1} prescribes for a vertex connected to an external line. Renormalisation families are formed inside each $H_K$ \cite[p.~666]{K1}.
\end{enumerate}
The only change from \cite[p.~666]{K1} is the cut criterion: King decides by vertex positions, (C3) by the clusters of the centres, which are summation variables of \cite[(3.36)]{K1} fixed in $\tau$.

\begin{lemma}[the cluster lemma]\label{lem:a02:CL}
Let $S$ be a prefix component of an ordering of a cluster graph $H_K$ with inside family $F_S$. \textup{(a$'$)} On the support of the localisation weights the centres of one vertex lie in one cluster. \textup{(a)} For $S'\in F_S$ and $H'\defeq(H\setminus S)\cup S'$: $H'$ carries the same tuple under the bijection of \cref{lem:a02:Lsharp}(c), hence the same clusters, cubes $\square_K,\square'_K,\square''_K$, points $x_0^K$, gauges $B_K$, fields $\tilde A_K$ and the same branch; every vertex and line of $S'$ lies in $K$; the legs of $S'$ are the legs of $S$, each with the same cut/uncut status; every line of $H'$ outside $S'$ has its status in $H$. The partition depends on no vertex position. \textup{(b)} On the supports a cut line $u$--$v$ has $|u-v|\ge R_0=3\sqrt d$, and a cut external line from $v$ to the block of a $y$-point has length $>R_0$. \textup{(c)} $\operatorname{diam}P_K\le c_0(|P_K|-1)$, every vertex of $K$ lies within $\frac23\sqrt d$ of $P_K$, and $|\square'_K|\le C(p+c_0N+M)^d\le C'p^d$, $N\defeq|P(\tau)|\le C(\nbar)$. \textup{(d)} The quotient graph (clusters as nodes, cut lines as edges) is connected, and for every spanning tree $T$ of it, $\dist(P(\tau))\le\sum_{\ell\in T}|\ell|+C(N)$, $|\ell|$ the distance between the ends of $\ell$.
\end{lemma}
\begin{proof}
(a$'$) Two centres $c,c'$ of a vertex $v$ have $v\in\supp\chi_{z_c}\cap\supp\chi_{z_{c'}}$, so $|z_c-z_{c'}|\le\frac43\sqrt d<c_0$ and they are adjacent in $\Gamma_\tau$. (a) The clusters are a function of $P(\tau)$, which is common by \cref{lem:a02:Lsharp}(c), and the cluster data are functions of $P_K$; the centres of $S'$ are those of $S$ as points, all in $K$ because $S$ is connected through uncut lines of $H_K$, so by (a$'$) every vertex of $S'$ is in $K$ and every line of $S'$ joins two vertices of $K$; a leg of $S'$ joins a vertex of $S'$ to the same outside end as the corresponding leg of $S$, in some cluster $K_w$, and its status is ``cut iff $K_w\ne K$'' in both; lines outside $S'$ are untouched; no step uses a position. (b) Let $z,z'$ be centres of $u,v$ in different clusters: $|z-z'|>c_0$ and $|u-z|,|v-z'|\le\frac23\sqrt d$. (c) A path in $\Gamma_\tau$ has steps $\le c_0$, and $p\ge1$. (d) Join the points inside each $P_K$ by a tree with steps $\le c_0$ and the clusters along $T$ through the centres of the endpoints of each $\ell$, at the cost $|\ell|+\frac43\sqrt d$ per edge.
\end{proof}
So both of King's phrases, ``treat it as an external line'' and ``the graph then factorizes into a product of localised subgraphs'' \cite[p.~666]{K1}, become exact, with the same factorisation for every member and every position of every family.

\begin{lemma}[regularised cut lines]\label{lem:a02:XL}
Fix a radial $\bar\psi\in C^\infty(\R^d)$, $0\le\bar\psi\le1$, $\bar\psi(w)=0$ for $|w|\le\frac12$, $\bar\psi(w)=1$ for $|w|\ge1$, and for a cut line $\ell$ between $u\in K$ and $v\in K'$ put $\ell^\flat(u,v)\defeq\bar\psi(u-v)\ell(u,v)$ (for an external line to a $y$-point, $\bar\psi(u-y)$), and $\ell^{\flat,\mathrm{ph}}\defeq U(B_K(\Gamma_{x_0^K,u}))^{-1}\ell^\flat U(B_{K'}(\Gamma_{x_0^{K'},v}))$ for the phase-dressed form of (C4). \textup{(a)} Replacing every cut line by $\ell^\flat$ leaves $E^{(j)}(H;\tau)$ unchanged, term by term: for every member of every family, every ordering and every position configuration, each evaluation pair $(u',v')$ at which a cut line enters, the lattice differences of the vertex functions at its ends included, has $|u'-v'|\ge R_0-4\cdot5^{-j}>1$, where $\bar\psi=1$. \textup{(b)} $\ell^{\flat,\mathrm{ph}}$ is $C^{1,\alpha}\otimes C^{1,\alpha}$ at unit scale in both arguments with decay $e^{-\delta_1|u-v|}$, at all positions and with no distance condition, the constant depending on $\nbar$, $d$, $L$, $\alpha$ only, times $(1+k-j)^5$ for a dressed line (the size of the coarse $\dm$ blob); this is \cite[(3.1)]{Ba3} in the form of a cut line with the endpoint localisations.
\end{lemma}
\begin{proof}
(a) The ends of a cut line are localised in the original sum: (C2) assigns to a cluster localised vertices, whose weights are present, block and $y$-points, which are points of $P(\tau)$, the vertices of one-centre $E_1$ pieces, none of whose lines is cut, and unlocalised coarse $\dm$ vertices, which lie strictly inside dressed lines and are never an end; so \cref{lem:a02:CL}(b) applies at every evaluation pair up to the lattice differences, each of which moves an end by at most $5^{-j}$. (b) For $|u-v|\ge\frac12$ the slice bounds \cite[Thm.~3.3, Props.~3.7--3.9]{K1} (\cref{prop:a02:FL}, \cref{prop:a02:S}) of the whole line, its derivatives and its mixed second differences are those of a unit-scale kernel with the stated decay, and the phases \cite[(3.46)]{K1} and $\bar\psi$ are smooth at unit scale; for $|u-v|<\frac12$ the kernel vanishes.
\end{proof}

\begin{theorem}[the Ward rows at every tuple]\label{thm:a02:WRall}
Let the localisation be by centres and the cut by clusters, with every cut line replaced by $\ell^\flat$, and let $S$ be a Ward-row prefix component of an ordering of a cluster graph $H_K$ at a fixed tuple $\tau$ (in particular in the branch $R>p$), with $A'$-line legs and scalar legs that may be cut lines. Then the sum over $F_S$ and $A_S$ is a sum of terms of the shapes of \cref{thm:a02:WR} with $\square\to\square_K$: \textup{(R0)} zero; \textup{(R1)} $\sum_x5^{-3j}c_a(x)\prod\partial^a\tilde\Phi(x)$, $|a|\le-D(S)$, $|c_a|\le C\rho_{j_0}^{D(S)+|a|}$, times $W^\circ$ and the frozen $\tilde A_K(x)$-legs, a cut-line leg counting as index $j$; \textup{(R2)} a credit $\rho^\theta$ on a line of $S$ with a debit of order $\le\theta$ on one leg (a cut-line leg costs $1$, by \cref{lem:a02:XL}(b)) or on a unit-scale weight, $D'(S)>0$, $D'(\gamma)\ge D(\gamma)$ for $\gamma\subsetneq S$. $C$ is uniform in $x$, $j_0$, $k$, $\eta$ and $\tau$ and depends on $\nbar$, $d$, $L$, $\alpha$ only, up to the dressed-line factor $(1+k-j)^5$ of \cref{lem:a02:XL}(b), absorbed in $(1+k)^{p_0}$ of \cref{cor:a02:DEGall}(b); the decay budget is $\nbar+4$ draws per line, and every evaluation point of a cut-line end is joined to a localised anchor of its own cluster by at most $C(\nbar)$ lines of rewritten units.
\end{theorem}
\begin{proof}
The proof of \cref{thm:a02:WR} is run inside $H_K$ with these changes. By \cref{lem:a02:CL}(a) every member of every family met in the nesting is reduced with the same cluster data, and cut lines never lie inside a unit: no ordered subgraph contains a cut line (it carries index $j$ and is external to both clusters), so a loop through a cut line is never a divergent unit. The frozen-weight identity of \cref{lem:a02:FW} holds per cluster, the legs that are cut lines entering through $\ell^{\flat,\mathrm{ph}}$, which by \cref{lem:a02:XL}(b) is a $C^{1,\alpha}$ leg at unit scale, so its debit in the Taylor and transport steps costs $1$; the identities of \cref{lem:a02:GC,lem:a02:REF} act only on $F^C_S$ after \cref{lem:a02:RS}, inside one cluster, and the cut lines are legs that enter only through the transport and Taylor steps after the identity, so the cut adds no identity and truncates no member line. The cut-line factor of an output term is bounded by $Ce^{-\delta_2d(K,K')}$ times transfer factors on unit lines: each evaluation point of a cut-line end is joined to a localised anchor by at most $C(\nbar)$ lines of its own cluster's rewritten units, the transfer takes one more decay draw per line, and the bound factorises over clusters. Nesting is as in \cref{thm:a02:WR}(4), with $\square_K$ in place of $\square$.
\end{proof}

\begin{corollary}[the degree count on every row; the bound in the branch $R>p$]\label{cor:a02:DEGall}
\textup{(a)} With the localisation by centres and the cut by clusters, every source-free class with $D\le0$ obeys \DEG\ at every tuple: the non-Ward rows by their member-wise identities (charge parity, the derivative at a line end, reflection, the bounded $\eta$-suppressed tadpole) or by the counterterm matching between partners of identical vertex structure, whose centre multisets coincide by \cref{lem:a02:Lsharp}(b) (the only non-Ward $D=-1$ class with a Taylor step is $\Sigma_1$, whose two centres sit at its two vertices, the attachment points of its scalar legs, so its non-reference centre is folded into the finer leg, \cite[(3.85)]{K1}, exactly as under King's rule); the Ward rows by \cref{thm:a02:WR,thm:a02:WRall}. \textup{(b)} For every graph whose divergent units are such classes, $|E^{(\theta)}(H_{\mathrm{ren}};\tau)|\le C\eta^{\rho_0}p^{n_1}(1+k)^{p_0}\exp[-\delta\dist(P(\tau))]$, $\theta\in\{k,k+n\}$, at every tuple, which is \cite[(3.38)]{K1} in the branch $R>p$; with \cite[(3.39)]{K1} from \cref{thm:a02:WR} in the branch $R\le p$ this gives King's per-graph estimate \cite[(3.42)]{K1}. \textup{(c)} At tuples with at least two clusters every connected component of every cluster graph carries a cut-line leg, so the vacuum members and the component-wise cancellation of \cite[p.~667]{K1} do not arise there, and $E(H;A^{(j)};\tau)$ and $E(H;0;\tau)$ are bounded separately; a closed scalar loop of internal lines of one cluster graph none of whose vertices carries another charged field vanishes when the total $q$-power on it is odd (no phase sits on the loop), which is the per-cluster form of charge parity.
\end{corollary}
\begin{proof}
(a) is the list of mechanisms of \cref{ssec:a02:ward} with the two theorems. (b) Per cluster, \cref{thm:a02:WRall} and (a) give the rewritten form with every stage of positive degree; the multiscale bound \cref{thm:a02:S3}(a) with the root held, the cut-line factors of \cref{lem:a02:XL} and \cref{lem:a02:CL}(d) (the tree distance of the tuple is controlled by the cut lines of a spanning tree) give the exponential in $\dist(P(\tau))$, and the volume factors $|\square'_K|\le C'p^d$ of \cref{lem:a02:CL}(c) give $p^{n_1}$. (c) A cluster graph is connected to the rest of $H$ through cut lines only, and $H$ is connected; on a closed loop with no charged external field the internal phases cancel and $\Tr q^{2r+1}=0$.
\end{proof}

For \cref{thm:a02:T} this supplies the $f=1$ input of (ii) on the Ward rows, $|c_a^1|\le C\rho_{j_0}^{D(S)+|a|}$, at every tuple of either branch; the transfer to the $f_D$ flow is (i)--(iii) there, the exact-zero and Ward mechanisms being identities for every charge-diagonal line kernel when partners carry the same lines, and \cref{lem:a02:Lsharp}(b) holding for the $f_D$ members verbatim since the weights do not enter the centre count.


\begin{theorem}[(B4), (B5): the curvature classes and the source-free classes in zone~I]\label{thm:a02:B4}\label{thm:a02:B5}
Let $H$ be a graph of King's expansion at step $j$ on $\CK$ at $h=0$ that is field-dependent, or contains a curvature vertex, and let it be localised in zone~I. Then \cite[(3.39)]{K1} holds for it, and \cite[(3.38)]{K1} with the extra factor $(1+k-j)^{p_0}$; summed over positions by \cref{lem:a02:SC}, the zone-I contribution to $|S^{(j)}_{N,k}-S'^{(j+1)}_{N,k+1}|$ is $\le C[5^{-\gamma'j}\eta^{-\beta'}+\eta^{\sigma'}]|\CJ|$, with $C$ depending on the curvature sources through $\max\norm{f_{\mu\nu}}_{C^{2,\alpha}}$ only.
\end{theorem}
\begin{proof}
By \cref{lem:a02:I} it suffices to treat $\tilde H^{\mathrm{free}}$ on the free $f_D$ flow, with the error of the form of \cite[(3.39)]{K1}. For the free flow run King's proof of \cite[Thm.~3.5]{K1} and of the step bound \cite[Prop.~3.6]{K1}: the census of ordered subgraph components (every component containing a source vertex has degree $\ge\frac12$, except the field-independent, source-linear family, which is absent at $h=0$ and vanishes for the curvature source at $A_j=\phi_j=0$ by charge parity, and the single-curvature-vertex tadpole, which vanishes by $\Tr q=0$); the ordered multiscale bound (\cref{thm:a02:S3}(a) with the degree count \DEG\ of \cref{thm:a02:T} for the source-free components); the one-factor-at-a-time replacement of \cite[pp.~663--665]{K1} with the line comparisons of \cref{prop:a02:FL} and, for the curvature vertex, \cref{lem:a02:SBP}(d), which carries the gain $\eps_k\le5^{-\gamma j}$ with a bounded summed vertex; and the identities \cite[(3.43)--(3.55)]{K1}, under which the curvature vertex is invariant (\cref{lem:a02:SBP}(b)). The curvature source enters the proof only through the sup and $(1,\alpha)$-norms of $w$ (\cref{lem:a02:SBP}(c)), uniformly in $\eps$. The free-boundary counterterm renormalises the free graphs exactly (\cref{lem:a02:I}).
\end{proof}

\subsection{Block vertices and the split degree}\label{ssec:a02:blocks}

King's block products \cite[(3.24)$\times$(3.26), (3.25)$\times$(3.26), (3.26)$\times$(3.26)]{K1} carry one or two summation points $x,x'$ in the block $B^j(y)$ and photon fields integrated along the staircases $\Gamma^{(j)}_{y,x}$. Following \cite[(1.13)--(1.15), (2.1)--(2.2), (2.12)--(2.14)]{Ba3} we read a graph $\tilde H$ with block products as the \emph{split graph} $G^{\mathrm{sp}}$: the points of the block products are vertices with measure $5^{-3j}$ and no explicit power of the spacing; a photon field on a staircase, contracted at slice $i$, enters through the staircase-integrated kernel $g^{(i)}(x,\cdot)\defeq G_{(i)}(\Gamma^{(i+1)}_{x_{i+1},x},\cdot)$, read as a scale-summed line of dimension $d_l\defeq2-3-\#\text{der}+c_l$, $c_l\in\{0,1,2\}$ the number of its staircase ends; the block kernel $\kappa_y(x,x')\defeq\mathbf 1[x,x'\in B^j(y)]$ (times $-a_j$ or $-a_j/2$) is a line of the single index $j$ and dimension $0$. The \emph{split degree} $\Dsp(c)\defeq\sum_{v\in c}w_v-3$ of a connected subgraph $c$ counts, at each vertex, $3$ plus the explicit powers of the spacing, minus $\frac12$ per undifferentiated and $\frac32$ per differentiated end of a line of $c$, plus $1$ per photon field of a point whose line lies in $c$; it is Ba{\l}aban's degree \cite[(2.1)--(2.2)]{Ba3}. (King's \cite[(3.30)]{K1} credits ``$+m$'' at a block vertex for all its $m$ photon fields; on subgraphs only the fields whose lines lie in the subgraph may be credited, as Ba{\l}aban's text says.)

\begin{lemma}[dictionary and telescoping]\label{lem:a02:Dsp}
(i) For a connected all-scale-summed subgraph $H$ of $G^{\mathrm{sp}}$ with components $c_1,\dots,c_r$ after removing the $\kappa$-lines, the degree of the multiscale lemma, $D_{MS}(H)\defeq\sum_{l\in H}d_l+\sum_{v\in H}a_v+3(|V(H)|-1)$, equals $\sum_i\Dsp(c_i)+3(r-1)$. (ii) For an ordering $l(1),\dots,l(m)$ with indices $i_1\le\dots\le i_m$,
\[
\prod_l\rho_{i_l}^{d_l}\prod_v\rho_{i(v)}^{a_v}\prod_{\text{tree lines}}\rho_{i_l}^{3}=\prod_{t=1}^m\rho_{i_t}^{D(H_t)-D(H_{t-1})}=\Big[\prod_{t<m}\Big(\frac{\rho_{i_t}}{\rho_{i_{t+1}}}\Big)^{D(H_t)}\Big]\rho_{i_m}^{D(H_m)},
\]
$i(v)$ the index of the first line at $v$. So a factor attached to a line enters exactly the stages containing that line, and a weight attached to a vertex enters from its first line on.
\end{lemma}
\begin{proof}
(i) Each end of a $c$-line at $v$ contributes $-\frac12$ or $-\frac32$ to $w_v$, so a line contributes $2-3-\#\text{der}$; a staircase end contributes $+1$ at its point, booked on the line as $c_l$; sum, and add the $\kappa$-lines ($d_\kappa=0$), which join the $r$ components. (ii) Adding $l(t)$ changes $D$ by $d_{l(t)}$, plus $a_v$ for each new endpoint, plus $3$ if $l(t)$ is a tree line.
\end{proof}

\begin{lemma}[line bounds for the split graph]\label{lem:a02:Bsp}
(i) $|g^{(i)}(x,b)|\le C\rho_ie^{-\delta_0|x-b|/\rho_i}$ (one staircase end; \cite[(3.64)]{K1}, \cite[(2.12)]{Ba3}); directly from \cite[(3.63)]{K1}, since $\Gamma^{(i+1)}_{x_{i+1},x}$ has at most $15\cdot5^{j}\rho_i$ bonds each weighted by $5^{-j}$ and lies within $5\sqrt3\rho_i$ of $x$. (ii) Two staircase ends: $C\rho_i^{2}$. (iii) Derivatives and H\"older differences in the non-staircase variable: \cite[(3.63), (3.65)]{K1} summed over the staircase. (iv) The finer model's slices $-1\le i$: the same. (v) Comparison across the pair: \cite[(3.74)]{K1} for the staircase line; its H\"older form from (3.74) and (3.75), the finer staircases subdividing the coarser ones except for the initial piece in the fine block, which costs $C5^{-j}\rho_i^{-1-\alpha}\le C5^{-\gamma j}\rho_i^{-\alpha-\gamma}$. (vi) $|\kappa^{(i)}(x,x')|\le e^{\delta\sqrt3}\delta_{ij}e^{-\delta|x-x'|}$. (vii) On $\CK$, (i)--(iii) hold with the cube slice bounds (\cref{prop:a02:S}(c), \cref{lem:a02:P}); the staircases lie inside the blocks, hence inside $\CJ$.
\end{lemma}

\begin{theorem}[the split degree governs the multiscale bound]\label{thm:a02:S3}
Let $\mathcal R$ be a renormalised family of King's expansion of order $\le\nbar$, written as a finite sum of graphs $G$ with the source vertices admitted, such that for each admitted ordering of the lines of $G^{\mathrm{sp}}$ every component of every ultraviolet prefix has current degree $D'>0$ after the rewriting (\DEG\ in the split count), the infrared lines obey \cite[(3.78)]{K1}, and the lines obey \cref{lem:a02:Bsp}. Let $\alpha_0$ be the minimum of these degrees. Then, on the torus and on $\CK$:
\begin{itemize}[nosep]
\item[(a)] the sum over orderings, compatible indices and positions of $|E^{(j)}(G)|$ is $\le C\eta^{\rho_0}p^{n_1}(1+k)^{p_0}\exp[-\delta\dist(\text{external vertices})]$, and the same holds with one element weakened by $\theta<\alpha_0$ (a line with $c_l\to c_l\rho_{i_l}^{-\theta}$, a vertex with an explicit power $a_v\ge\theta$, or an external field);
\item[(b)] $|E'^{(j+1)}(G)-E^{(j)}(G)|\le C\eta^{\rho'_0}5^{-\gamma j}\exp[-\delta\dist]$ for $0<\gamma<\alpha_0$;
\item[(c)] every rewritten subgraph is source-free, point-free and a component of a prefix of $G^{\mathrm{sp}}$, and $D'(\gamma)\ge\Dsp(\gamma)$ for every component $\gamma$ of every prefix;
\item[(d)] the exponent of every stage of the ordered bound is $D'(H_t)$ in the split count; no per-vertex credit for staircase fields enters.
\end{itemize}
\end{theorem}
\begin{proof}
(a) Apply the multiscale lemma of \cite[(3.66)--(3.70)]{K1} (a line-by-line induction over the ordering, each step shrinking one line and summing one endpoint, the index sums geometric by the positive degree) to the connected graph $G^{\mathrm{sp}}$ with the lines and bounds of \cref{lem:a02:Bsp}. Its degree is $\Dsp$ by \cref{lem:a02:Dsp}(i); a subgraph containing a $\kappa$-line has $D_{MS}=\sum_i\Dsp(c_i)+3(r-1)$ with each $c_i$ containing a point, hence $\ge1$ by \cref{lem:a02:BLK}. Infrared stage: points lie in no counterterm unit, so no infrared line ends at a point; $\kappa$ has index $j\le$ every infrared index, and shrinking it (put it first among the ties) transfers only external lines at $x'$ to $x$ and gives $\sum_{x'}5^{-3j}\kappa\le1$, shifting every later prefix degree by the same constant; so the infrared differences of \cite[(3.80)]{K1} are unchanged. Fixed lines, remainders and the position sum as in \cite[pp.~663--664]{K1}; the weakened element by the same induction with $D'-\theta\ge\alpha_0-\theta>0$. (b) \cite[pp.~663--665]{K1} verbatim on $G^{\mathrm{sp}}$: the finer slices give $5^{-\gamma j}$ through the finest sum; otherwise replace one factor at a time (external lines by (3.71), sources by sampling, the phase by (3.72), plain lines by (3.73), staircase lines by (3.74), H\"older lines by (3.75) and \cref{lem:a02:Bsp}(v)), each error being the graph with one element weakened by $\gamma$; then replace the point sums, $\sum_{x'\in B^1(x)}5^{-3(j+1)}=5^{-3j}$ exactly, and $\kappa'=\kappa$ exactly since the unit blocks of both members coincide. (c) A rewritten subgraph is divergent, so by \cref{lem:a02:BLK} it contains no point and is one of Ba{\l}aban's divergent classes, which contain no block vertex \cite[pp.~428--431]{Ba3}; debits land on legs of $S$, and if a leg lies in $\gamma$ then $S\subset\gamma$, so the net change per component is $\ge0$. (d) is \cref{lem:a02:Dsp}(ii). On $\CK$ use \cref{lem:a02:Bsp}(vii) and bound position sums over $\CJ$ by sums over $\Z^3$.
\end{proof}

\begin{lemma}[block-vertex components are convergent]\label{lem:a02:BLK}
Let $\gamma$ be a connected source-free component of King's expansion at zero background, on $\CK$ or on the torus, containing a point of a block product and properly contained in a connected graph. Then $\Dsp(\gamma)\ge1$, or $\gamma$ vanishes identically; $\Dsp(\gamma)\ge2$ if $\gamma$ contains a second point or a source vertex.
\end{lemma}
\begin{proof}
A proper connected subgraph has $E_\gamma\ge1$ external legs. Every King vertex has an even number of scalar legs and each internal scalar line uses two ends, so $E_\gamma$ has the parity of the number of external photon legs. With $D(\gamma)=\sum_vu_v-3+E_\gamma/2$ \cite[(3.30)--(3.31)]{K1}, a point with a photon field whose line lies in $\gamma$ has $w\ge\frac52$ and every other vertex $w\ge0$: if $E_\gamma\ge2$, $\Dsp\ge\frac52-3+1=\frac12$, hence $\ge1$ by integrality ($\Dsp$ is an integer, since $a_v$ and the number of differentiated internal ends are integers and $\sum_vn_v=2\#\text{lines}+E_\gamma$). If $E_\gamma=1$ the leg is a photon; at zero background every scalar line is component-diagonal and every vertex carries $q^{m_v}$, so each closed charged loop carries $\Tr q^{n_c}$ with $\sum_cn_c$ odd, and $\Tr q^{\mathrm{odd}}=0$. A second point or a source vertex adds $\ge\frac32$ (curvature) or $2$ (point, scalar source).
\end{proof}

\begin{lemma}[the position sum]\label{lem:a02:SC}
Let a family of localised graphs on $\CK$ satisfy, for each localisation, the root-summed bound $\sum_{\text{other positions}}|E^{(j)}(H;\{y\},\{z\})|e^{\delta'\dist}\le C\eta^{\rho_0}p^{n_1}(1+k)^{p_0}\Lambda(y_1)$ with the root $y_1$ fixed and a root weight $\Lambda$ with $\sum_{y_1\in Q}5^{-3j}\Lambda(y_1)\le C$ for every unit cube $Q$. Then the sum over all localisations of the field- and source-dependent part is $\le C\eta^{\rho_0}p^{n_1}(1+k)^{p_0}\eta^{-3}|\CJ|$ for the $R\le p$ branch, and $\le C\eta^{3+\sigma}\eta^{-3}|\CJ|$ for the $R>p$ branch, where the factor $e^{-(\delta-\delta')p}$ is $\le C\eta^{3+\sigma}$ because $p(\eta)$ grows faster than every logarithm. On the diagonal sequence $j=j(k)$ of \cref{thm:k01:CA} the powers $p^{n_1}(1+k)^{p_0}$ are $\le C_\epsilon5^{\epsilon j}$, so a bound $5^{-\gamma j}\times$(this) is $\le C5^{-\gamma'j}\eta^{-\beta'}|\CJ|$ for any $\gamma'<\gamma$ and $\beta'\defeq\max(3-\rho_0,0)$.
\end{lemma}
\begin{proof}
This is \cite[(3.40)--(3.42)]{K1} with the position sum over $\eta^{-3}|\CJ|$ unit sites and the root weight summed by its hypothesis; the counterterm units' vertices off $T_\eta$ are internal to a constant and do not enter.
\end{proof}

\subsection{Zone B: the scalar face remainders}\label{ssec:a02:Bsc}

In zone~B the images of the lines across the faces on $F$ are at distance $2t(x)$, not $\ge4p$, so \cref{lem:a02:I} fails, and the free-boundary counterterm renormalises the Dirichlet self-energies only up to remainders \cite[pp.~331--332]{K2}. We write the cube rewriting as the free one plus remainders, and bound the remainders by face-layer weights.

\begin{lemma}[the rewriting steps are identities on the cube]\label{lem:a02:SBPC}
(a) Taylor transport along axis paths between live sites of $\Omega$, H\"older splits, the resolvent splitting $G=C^\xi+G(C^{\xi,-1}-G^{-1})C^\xi$ with $C^{\xi}\defeq(-\Delta^\Omega+\xi^{-2})^{-1}$, rescaling and resummation over orderings are pointwise identities on $\Omega$. (b) On a bond with one end on a Dirichlet plane the hopping term is $|\phi(b_\pm)|^2$, independent of $A_b$; the vertices on such bonds sum to zero and are dropped, so every vertex sits on bonds with both ends live. (c) For $f$ vanishing off the live sites and $w$ supported on live bonds, $\sum_b(\nabla_bf)w(b)=-\sum_xf(x)(\nabla^*w)(x)$ exactly, with no boundary term: this is \cite[(2.8)]{Ba3} and \cite[(3.16)]{K2} on the cube.
\end{lemma}
\begin{proof}
Finite sums on $5^{-j}\Z^3$; (b) uses $|U(A_b)\phi(b_+)-0|^2=|\phi(b_+)|^2$.
\end{proof}

\begin{definition}[cube rewriting, remainder vertices]\label{def:a02:RC}
Apply the steps of the free rewriting to a zone-B family with cube lines and live-site positions. At each evaluation of a resummed coefficient $c_S(x)$ (a \emph{C-site} $S$, rooted at $x$) write $c^C_S(x)=c^F_S(x)+\Delta_S(x)$, $c^F_S$ the value with free $f_D$ lines and positions over $5^{-j}\Z^3$; the \emph{remainder} $\Delta_S(x)$ is attached at $x$ as a generalised \emph{remainder vertex} $v_S$ with the legs of the evaluated term, and every graph carrying a remainder is \emph{frozen}: no later step acts on a stage containing a remainder vertex. The \emph{moment degree} $D'(S)$ is $D(S)$ for a zeroth moment and $D(S)+1$ for a first moment.
\end{definition}

\begin{lemma}[comparison decomposition]\label{lem:a02:CMP}
Every zone-B renormalised family $\mathcal R$ of order $\le\nbar$ on $\CK$ is a finite sum $\sum G+\sum G'$ of \emph{bulk-type} graphs (the free rewriting's outputs with cube lines and live positions, free coefficients at the C-sites; they obey \DEG\ of \cref{thm:a02:T}, degrees being combinatorial, and their lines obey the cube slice bounds) and \emph{remainder} graphs (each containing a remainder vertex, with every stage containing it unrewritten). The free-boundary counterterm units are the partners of the bulk-type part exactly, as in zone~I.
\end{lemma}
\begin{proof}
Induction along the rewriting: \cref{lem:a02:SBPC}, the identity $c^C=c^F+\Delta$, and the fact that an inner remainder is never inside a collapsed coefficient (its graph is frozen before the outer evaluation).
\end{proof}

\begin{theorem}[census of the scalar remainders]\label{thm:a02:CEN}
Let $S$ be a C-site of a zone-B family other than the vacuum polarisation sites of \cref{ssec:a02:BVP} and the field-free vacuum components. Then (i) $S$ contains no block point and no source vertex; (ii) $D'(S)\in\{-1,0\}$, $S$ being the zeroth moment of a scalar self-energy of order $2$ or $4$, the first moment of the order-$2$ one-photon exchange, a one-leg-differentiated two-point class, or a $\phi\phi A$ class; (iii) $D'(S)=-1$ iff $S$ is the zeroth moment of an order-$2$ scalar two-point class with both legs undifferentiated (the one-photon exchange after \cite[(3.16)]{K2}, the $\lambda$-tadpole, the seagull tadpole); then $v_S$ has exactly two scalar legs, both undifferentiated at $v_S$; (iv) every connected non-vacuum source-free stage that is not charge-odd has $D\ge-1$. Moreover, every connected stage $H$ of a non-vanishing member that contains a line at a scalar remainder vertex and has $E_H\ge1$ satisfies $D(H)\ge0$, computed with $u_{v_S}=D'(S)+3-E_S/2-\epsilon_v$ ($\epsilon_v=1$ if a first-moment derivative sits on a leg internal to $H$), and $D(H)\ge1$ if $H$ contains a block point.
\end{theorem}
\begin{proof}
(i) A block point gives $\Dsp\ge1$ (\cref{lem:a02:BLK}), so no evaluation acts there; sources by \cref{lem:a02:Z}. (ii) is Ba{\l}aban's list of evaluations \cite[pp.~429--445]{Ba3} for the point-free vertex set, which is King's \cite[(3.18), (3.19), (3.22), (3.23)]{K1}, with the degrees from \cite[(3.30)--(3.31)]{K1}: for two scalar legs $D=\sum u-2$, the one-photon exchange (two current vertices with internal differentiated ends, $u=\frac12$ each) and the tadpoles ($u=1$) have $D=-1$, order $4$ has $D=0$, order $\ge6$ is convergent; a first moment adds $1$; a derivative on an external leg does not count, so the one-leg-differentiated classes have $D=0$; for $\phi\phi A$ legs $D=\sum u-\frac32\ge0$, and $D$ is an integer (the parity count below), so the next value is $1$. (iii) by (ii): the current vertex's derivative is on the internal line, the transported leg is undifferentiated. (iv) $E\ge5$: $D>0$ \cite[(2.17)]{Ba3}; $E=4$: $D\ge0$; $E=3$: $D\ge0$; $E=2$: $D\ge-1$; a scalar one-leg graph needs a point, a photon one-leg graph is charge-odd. \emph{The stage bound.} $v_S$ carries $a_v=D'(S)$ (\cref{lem:a02:K18} books $\rho_{j_0}^{D'}$ at $j(v)=j_0$), $n_v=E_S$ legs and $\delta_v=\epsilon_v$, so $u_v$ is as stated and $n_v\equiv2u_v\pmod2$: $D(H)$ is an integer. For a zeroth moment $u_v=\sum_{w\in S}u_w$, so $D(H)=D(\hat H)$ with $v$ re-expanded; a first moment adds $1-\epsilon_v\ge0$. Units met by $H$ through a counterterm line are replaced by point vertices of $u=2$, each whole order-$2$ unit costing $\ge-1$ and every other $\ge0$; with $k_1$, $k_2$ their numbers, $D(H)\ge D'(S)-E_S/2-\epsilon_v+k_1+2k_2+\sum_{\mathrm{std}}u+E_{H'}/2$. If $k_2\ge1$, $D(H)\ge0$; if $k_2=0<k_1$, $E_{H'}=E_H\ge1$ gives $D(H)\ge-\frac12$, hence $\ge0$. If $k_1=k_2=0$ put $X\defeq\sum_{w\ne v}u_w+E_H/2$, so $D(H)=D'(S)-E_S/2-\epsilon_v+X$ with $X-E_S/2\in\Z$. For $E_S=2$: $X\ge1$ (a self-loop at $v$ alone leaves $E_H=0$); $X=1$ forces $E_H=1$ and $H=\{v,w\}$ with $w$ a current vertex whose two scalar ends close both legs of $v$, a scalar cycle $v$--$w$--$v$ carrying $\Tr(M_vq)=0$ ($M_v\propto1$), so the member vanishes; hence $X\ge2$ and $D(H)\ge D'(S)-\epsilon_v+1\ge0$. For $E_S=3$: $X\in\frac12+\Z$, $X=\frac12$ forces a self-loop closing the two scalar legs of $v$, whose cycle carries $\Tr q^{\mathrm{odd}}=0$; hence $X\ge\frac32$ and $D(H)\ge D'(S)=0$. Block points: $D(H)\ge D(\hat H)\ge1$ by \cref{lem:a02:BLK}.
\end{proof}

\begin{lemma}[geometry of a remainder]\label{lem:a02:GEO}
Let a configuration of a C-site $S$ rooted at a live site $x$ contain a boundary element: an image line $(u,gv)$, $g\ne1$ in the reflection group of $\Omega$, or a free-side internal position $y$ that is not live. Then the lengths $l_l\defeq|u_l-g_lv_l|$ of the lines on a path from $x$ to the element sum to $\ge t(x)/2$; if $g$ involves a face of $\Omega$ not on $F$, then $l_l\ge2p$.
\end{lemma}
\begin{proof}
$|u-gv|\ge\dist(u,\partial\Omega)+\dist(v,\partial\Omega)$ and $|a-b|\le|a-gb|$ for every path line, so $t(x)\le|x-u|+t(u)\le\sum_{\text{path}}l+l_l$, up to the photon half-offset $t_N(u)\ge t(u)/2$; an outside vertex gives $t(x)\le|x-y|$; the vertices of $S$ lie in $\square$, at distance $\ge2p$ from the faces not on $F$.
\end{proof}

\begin{lemma}[layer weight of a remainder vertex]\label{lem:a02:K18}
Let $S$ be a scalar C-site, $j_0$ the lowest index of its legs, $\rho\defeq\rho_{j_0}$, $D'\defeq D'(S)$, $x$ live, $t\defeq t(x)$. Then
\[
|\Delta_S(x)|\le C\rho^{D'}\big[w_{D'}(t/\rho)+e^{-\delta p/\rho}\big],\qquad w_{-1}(s)\defeq s^{-1}e^{-\delta s},\quad w_0(s)\defeq(1+\ln_+(1/s))e^{-\delta s},
\]
and across the pair $|\Delta'_S(x)-\Delta_S(x)|\le C5^{-\gamma j}\rho^{D'-\gamma}[(\rho/t)^\gamma w_{D'}(t/\rho)+e^{-\delta p/\rho}]$ for $0<\gamma<\alpha_0$. $C,\delta$ depend on $\nbar$, $\alpha$, $\alpha_0$ and the line constants only.
\end{lemma}
\begin{proof}
Expand the internal lines of $S$ into slices of index $\le j_0$ and fix an ordering; expand $\Delta_S$ multilinearly, each term having a boundary element (\cref{lem:a02:CMP}); image lines are free kernels at $(u,gv)$ with $|K^{\mathrm{free}}(u,gv)|\le C\rho_i^{d_l}e^{-\delta_1l_l/\rho_i}$, and the image sums converge since the sides of $\Omega$ are $\ge1\ge\rho$. Run the multiscale induction (\cref{thm:a02:S3}(a)) with the root held, splitting each line's kept half-exponential into quarters (external distance, transfer, image count, \cref{lem:a02:GEO}); since $\rho_{i_l}\le\rho_{i_m}$ for the top internal index $i_m$, \cref{lem:a02:GEO} gives $\prod_le^{-(\delta_1/8)l_l/\rho_{i_l}}\le e^{-(\delta_1/16)t/\rho_{i_m}}$, or $e^{-\delta_1p/(4\rho_{i_m})}$ for an off-$F$ image. The proper stages of $S$ have $D\ge\alpha_0>0$ (\DEG: ``the only divergent graphs can be the whole graphs'', \cite[p.~445]{Ba3}), so their index sums are geometric, and the telescoped product (\cref{lem:a02:Dsp}(ii)) leaves $\rho_{i_m}^{D'}$ (a first-moment factor $|x'-x|$, distributed along a chain with $|a-b|\le|a-gb|$, adds $+1$). Hence $|\Delta_S(x)|\le C\sum_{i_m\le j_0}\rho_{i_m}^{D'}e^{-\delta't/\rho_{i_m}}+(\text{off-}F)$, which is $\le Ce^{-\delta t/\rho}(1+\ln_+(\rho/t))$ for $D'=0$ and $\le Ct^{-1}e^{-\delta t/\rho}$ for $D'=-1$; the off-$F$ part is $\le C_Ne^{-\delta p/2}$ for every $N$. The pair difference: replace one line at a time by its difference (\cite[Props.~3.8--3.9]{K1} on $\Omega$, via \cref{prop:a02:S}(c)), which weakens it by $\rho_{i_l}^{-\gamma}$ with the prefactor $5^{-\gamma j}$; the finer slices $i<0$ give $5^{-\gamma j}$ through the finest stage ($D(H_1)\ge\alpha_0>\gamma$), a top-only finer line (the $\lambda$-tadpole) giving $\sum_{i<0}\rho_i^{-1}e^{-\delta't/\rho_i}\le Ct^{-1}(5^{-j}/t)^\gamma$; the point-sum replacement is exact since both members have the same planes.
\end{proof}

\begin{lemma}[the borrowed H\"older factor]\label{lem:a02:H}
Let $x$ be live, $t=t(x)$, $0<\alpha<1$, $m_\alpha(x;\rho)\defeq\min(1,t/\rho)^\alpha$. (i) For a cube slice with any far-end decoration $O_z$ of order $\mathsf q\in\{0,1,1+\alpha'\}$: $|O_zG^\Omega_{(i)}(x,z)|\le C\rho_i^{-1-\mathsf q}m_\alpha(x;\rho_i)e^{-\delta|x-z|/\rho_i}$. (ii) $|\phi^{(j)}(x)|\le Cm_1(x;1)\norm{\phi_j}_\infty$ and $|a_jG_jQ_j^*(x,y)|\le Cm_\alpha(x;1)e^{-\delta|x-y|}$. (iii) $|\delta G_j(\square',A^{(j)})(x,y)|\le Cm_\alpha(x;1)e^{-\delta_0\dist(\{x,y\},\partial\Omega\setminus F)}$. (iv) If the line also carries a pair difference (bound $5^{-\gamma_1j}\rho^{-\gamma_1}\times$ its bound), its difference obeys the difference bound with $\gamma=\gamma_1/2$ times $m_{\alpha/2}$.
\end{lemma}
\begin{proof}
(i) The mixed clause: each slice factorises as $G_{(i)}=a_i^2\rho_i^{-4}(G_iQ_i^*)C^{(i)}(Q_iG_i)$ \cite[(2.17), (3.60)]{K1} with box operators on $\Omega$, so the $x$-dependence sits in the left factor and the $z$-dependence in the right one, and $|O_zG_{(i)}(x,z)-O_zG_{(i)}(x',z)|\le C|x-x'|^\alpha\rho_i^{-1-\alpha-\mathsf q}e^{-\delta\dist(\{x,x'\},z)/\rho_i}$ by the H\"older clause of the left factor and the decorated clause of the right one; with $x'=x_P$ the foot of $x$ on the nearest Dirichlet plane, where $G_{(i)}(x_P,\cdot)=0$, this is the claim for $t\le\rho_i$, and for $t>\rho_i$ use the single-point bound. This is King's \cite[(3.19)]{K2} for every leg type. (ii) $\phi^{(j)}$ vanishes on the Dirichlet planes and $|D_A\phi^{(j)}|\le C\norm{\phi_j}$ over face bonds \cite[(5.2)]{K2}; telescope along the $t5^{j}$ bonds from $x$ to $x_P$. (iii) Both kernels vanish on $F\cap\partial\Omega$, and \cite[(5.3)]{K2} gives the derivative and H\"older clauses of the difference with decay from the faces not on $F$. (iv) $\min(a,b)\le a^{1/2}b^{1/2}$, both members' slices vanishing on the same planes.
\end{proof}

\begin{lemma}[charging rule and the weight after borrowing]\label{lem:a02:CH}
In a non-vanishing member of a remainder graph one can choose for each $D'=-1$ remainder vertex $v$ one of its two scalar leg ends at $v$ so that every line is charged at most once. Fix $0<\beta<\alpha/2$ and move the factor $m_\beta(v;\rho_l)$ of the charged leg $l$ (\cref{lem:a02:H}) from the line to $v_S$. Then the line keeps its standard bound, and $v_S$ carries $\omega'_S(x)\defeq\rho|\Delta_S(x)|m_\beta(x;\rho_l)\le C[(\rho/t)^{1-\beta}e^{-\delta t/\rho}+e^{-\delta p/(2\rho)}]$, a weight of class $\lay\rho0{1-\beta}$ (\cref{def:a02:lay}) plus an off-$F$ term of any degree; across the pair the difference is $5^{-\gamma j}\rho^{-\gamma}$ times a weight of class $\lay\rho0{1-\beta+\gamma}$, $0<\gamma<\beta$.
\end{lemma}
\begin{proof}
Let $\mathcal R$ be the multigraph on the $D'=-1$ remainder vertices with the member lines joining two of them as edges. Each node has degree $\le2$, so $\mathcal R$ is a union of paths and cycles; a cycle uses both legs of each node and is a whole component with $E=0$, a field-free vacuum component, which is dropped. Orient each path and charge each node to its left leg end; the leftmost node's left leg is a line to a vertex outside $\mathcal R$ or an external field, and each edge is charged once, at its right end. The weight: $m_\beta(x;\rho_l)\le m_\beta(x;\rho)$ for $\rho_l\ge\rho$, and \cref{lem:a02:K18}. Without the borrowing the $D'=-1$ layer sum diverges like $\ln(\rho5^{j})$ per unit face area; with any $\beta>0$ it is bounded. This is King's ``extra degree of convergence \ldots borrowed from a line external to $H$'' \cite[p.~332]{K2}, with the rule that no line gives up two factors.
\end{proof}

\begin{theorem}[(B5$'$), scalar part]\label{thm:a02:Bsc}
Let $\mathcal R$ be a renormalised family of order $\le\nbar$ of King's expansion of $P^{(j),1}$ on $\CK$ at zero background, localised in zone~B, with at least one external field or at least one source vertex, rewritten by \cref{def:a02:RC} (its vacuum-polarisation prefix components by \cref{thm:a02:BVP}). Fix $0<\alpha<1$, $0<\beta<\alpha/2$ and $0<\gamma<\min(\alpha_0^B,\beta)$, $\alpha_0^B$ the minimum positive stage degree over the finitely many graphs and orderings. Then:
\begin{itemize}[nosep]
\item[(a)] with one vertex $x_0$ held (or none), the sum over orderings, compatible indices and positions of $|E^{(j)}(G)|$ over the members of $\mathcal R$ is $\le C\eta^{\rho_0}p^{n_1}(1+k)^{p_0}\exp[-\delta\dist(V_*)]\Lambda(x_0)$, $V_*$ the vertices with external legs, the held vertex and the source vertices, with $\Lambda(x_0)\le C(1+\ln_+(1/t(x_0)))^{\nbar+1}t(x_0)^{-(1-\beta)}e^{-\delta_\Lambda t(x_0)}$, $\Lambda=1$ if no vertex is held, $\Lambda(x_0)\le C$ if $t(x_0)\ge1$ (in particular at every source vertex), and $\sum_{x_0\in Q}5^{-3j}\Lambda(x_0)\le C$ for every unit cube $Q$;
\item[(b)] across the pair, $|E'^{(j+1)}-E^{(j)}|\le C5^{-\gamma j}\times$(the bound (a) with $t^{-(1-\beta)}$ replaced by $t^{-(1-\beta+\gamma)}$);
\item[(c)] summed over the zone-B localisations (\cref{lem:a02:SC}), the zone-B part of the source-free field-dependent effective-action difference is $\le C5^{-\gamma j}\eta^{\rho_0}p^{n_1}(1+k)^{p_0}|\partial T^{(j)}|\le C'5^{-\gamma'j}\eta^{-\beta'}|\CJ|$ on the diagonal sequence, for every $\gamma'<\gamma$ and $\beta'=\max(2-\rho_0,0)$; for source hosts the position sum is the host's, with the root at the source and the factor $e^{-\delta R/2}$, $R\ge\eta^{-1}/2$.
\end{itemize}
The constants depend on the point of King's range, $\nbar$, $\alpha,\beta,\gamma$ and the line constants, not on $j,k,\eta,J,x_0$.
\end{theorem}
\begin{proof}
\Cref{lem:a02:CMP} splits $\mathcal R$ into bulk-type and remainder graphs. \emph{Bulk-type graphs} are generalised graphs with \DEG\ (\cref{thm:a02:T}; with sources, the host's own census as in \cref{thm:a02:B4}), lines obeying the cube slice bounds, vertices bounded by sup norms (sources by $\norm{w}_\infty$, \cref{lem:a02:SBP}(c)); \cref{thm:a02:S3}(a) holds on $\CK$ verbatim (the induction uses no periodicity), so (a) holds with $\Lambda\equiv1$, and (b) is \cref{thm:a02:S3}(b) with the cube line differences. \emph{Remainder graphs} satisfy the hypotheses of \cref{thm:a02:OFB}: lines obey the standard bounds (a charged leg keeps its bound, \cref{lem:a02:CH}); the layered vertices are the $D'=0$ remainder vertices, of class $\lay{\rho_{j_0}}10$ (\cref{lem:a02:K18}), the $D'=-1$ ones after borrowing, of class $\lay{\rho_{j_0}}0{1-\beta}$ (\cref{lem:a02:CH}), and the vacuum-polarisation face pieces (\cref{thm:a02:BVP}(b)), each of width $\rho_v\le\rho_{j(v)}$ since $j(v)=j_0$, with off-$F$ terms of any prescribed degree. \DEGB: every connected component of every admitted stage has $D\ge0$, and $D=0$ only with a layered vertex: stages containing a scalar remainder vertex by \cref{thm:a02:CEN}; stages containing a vacuum-polarisation piece by \cref{lem:a02:CENpp}; stages containing a block point by \cref{lem:a02:BLK}; stages containing a source vertex and a layered vertex (\emph{mixed}) by \cref{lem:a02:SRC}; all other stages are bulk-type stages of the frozen graph with $D>0$ by \DEG. \Cref{thm:a02:OFB}(i) gives (a), with $\theta_{\max}=1-\beta<1$ so that the root weight is summable, and $\Lambda(x_0)\le C$ at $t(x_0)\ge1$. \Cref{thm:a02:OFB}(ii) gives (b): its pair hypotheses are the line differences (\cref{prop:a02:S}(c), \cref{lem:a02:P}), the coefficient differences (\cref{lem:a02:K18}, \cref{thm:a02:BVP}(c)), the layered weights in rate form (\cref{lem:a02:CH}, class $\theta+\gamma<1$), the sources by sampling and \cref{lem:a02:SBP}(d), and $\gamma<\min(\alpha_0^B,\beta)$. The infrared stage is unchanged: infrared lines lie inside free counterterm units, which contain no remainder, no source and no layered vertex. Field-independent members with a source vertex (curvature classes with $n_g\ge2$ even): $\mathcal R$ of \cref{lem:a02:CH} has no cycle there, because a cycle would be a whole component made of remainder vertices, and the member is connected and contains a source vertex, which is not a remainder vertex. (c): \cref{thm:a02:BVP}(d) verbatim, with $\sum_{x_0}5^{-3j}\Lambda(x_0)\le C$ per unit face cell (\cref{lem:a02:L}(b) at width $1$) and $|\partial T^{(j)}|\le C\eta^{-2}|\CJ|$.
\end{proof}

\begin{lemma}[stages containing a source vertex and a layered vertex]\label{lem:a02:SRC}
Let $H$ be a connected component of an ultraviolet stage of a non-vanishing remainder member, containing a source vertex $s$ and a layered vertex $v$. Then (a) in the induction of \cref{thm:a02:OFB}, reserving half of every ultraviolet line's decay and half of every layered weight's exponential, the reserved factors give $\le C_N\rho_{i_{t_*}}^N$ for every $N$ at the first step $t_*$ at which such a component appears, raising the exponent of every later ultraviolet step by $N$; (b) $D(H)\ge0$; (c) no C-site contains a source vertex.
\end{lemma}
\begin{proof}
(a) A path $P$ of lines of the stage joins $s$ to $x_v$, each of index $\le i_{t_*}$; successive vertices on $P$ are attached within one fine spacing, so $\sum_{l\in P}|l|\ge|s-x_v|-\nbar5^{-j}$, and with $t(s)\ge\eta^{-1}/2\ge4$ (\cref{lem:a02:Z}(iii)) and $t$ $1$-Lipschitz, $\sum_P|l|+t(x_v)\ge t(s)-\nbar5^{-j}\ge1$; also $\rho_v\le\rho_{i_{t_*}}$. Hence the reserved factors are $\le\exp[-c_1/\rho_{i_{t_*}}]\le C_N\rho_{i_{t_*}}^N$, and $\rho_{i_{t_*}}^N\le\prod_{t\ge t_*}(\rho_{i_t}/\rho_{i_{t+1}})^N$. (b) Re-expand the generalised vertices of $H$: $D(H)\ge\Dsp(\tilde H)$ (zeroth moments give equality, first moments and the bulk rewriting never lower a containing degree, \cref{thm:a02:S3}(c)); replace the units met through a counterterm line by points of $u=2$ as in \cref{thm:a02:CEN}. If some unit is met, $\Dsp\ge u_s-2$; otherwise $\tilde H$ contains $s$ and the vertex set of the layered vertex's class, whose standard vertices give $\sum u\ge1$, so again $\Dsp\ge u_s-2$. With $u_s=2$ (scalar source) or $\frac52$ (summed curvature vertex) this is $\ge0$. (c) A connected subgraph with a source vertex has $D\ge\frac12$ except the field-independent source-linear family (absent at $h=0$; for the curvature source it vanishes by charge parity) and the single-vertex tadpole ($\Tr q=0$): \cref{thm:a02:B4}'s census.
\end{proof}

\subsection{Zone B: the vacuum-polarisation face pieces and the ordered bound with layered vertices}\label{ssec:a02:BVP}

\begin{definition}[layered weights]\label{def:a02:lay}
A weight $\omega$ on $\Omega$ is \emph{layered} of class $\lay\rho m\theta$ ($5^{-j}\le\rho\le1$, $m\le\nbar$, $0\le\theta<1$, $\delta_b>0$) if $|\omega(x)|\le(1+\ln_+(\rho/t(x)))^m(\rho/t(x))^\theta e^{-\delta_bt(x)/\rho}$ for every site. A vertex is \emph{layered} if its vertex function is a bounded function times $\omega_v(x_v)$ with $\omega_v\in\lay{\rho_v}m\theta$ and $\rho_v\le\rho_{j(v)}$, $j(v)$ the lowest index of the lines at $v$.
\end{definition}

\begin{lemma}[integrality and the photon two-point census]\label{lem:a02:CENpp}
(i) For a connected $\gamma$, $D(\gamma)=\sum_vu_v-3+E_\gamma/2$ is an integer. (ii) Let $\gamma$ be connected, source-free, with exactly two external legs, both photon. Then $D(\gamma)\ge-1$, with equality iff $\gamma$ is the one-loop bubble (two current vertices joined by two scalar lines) or a tadpole (one seagull vertex with its scalar ends joined); all others have $D\ge0$. (iii) Let $m$ be a non-vanishing member, $S'$ a connected photon two-point prefix component and $H$ a connected component of a stage with $S'\subsetneq H$, $E_H\ge1$. Then $D(H)\ge0$, and $D(H)\ge D(S')+1$ if $D(S')=-1$.
\end{lemma}
\begin{proof}
(i) $a_v$ and the number of differentiated internal ends are integers, so $n_v\equiv2u_v\pmod2$, and $\sum_vn_v=2\#\text{lines}+E_\gamma$. (ii) $D=\sum u-2$; the two photon ends sit at vertices of total $u\ge1$, with equality only for the bubble and the tadpoles (block vertices have $u\ge\frac52$, and a $u=0$ merged counterterm vertex drags a $u\ge\frac12$ vertex of its unit into $\gamma$). (iii) $\sum_{v\in S'}u_v=D(S')+2$; $H$ contains a leg of $S'$ and its other end $w$, $u_w\ge\frac12$, and $D(H)=D(S')-1+(\Sigma'+E_H/2)$ with $\Sigma'\defeq\sum_{v\in H\setminus S'}u_v$, $\Sigma'+E_H/2$ an integer $\ge1$. If $D(S')=-1$ and $D(H)<0$ then $\Sigma'=\frac12$, $E_H=1$, $H\setminus S'=\{w\}$ a current vertex whose differentiated end closes a loop at $w$ with one photon vertex, $\Tr[V^{\mathrm{hop}}_bG_{(i)}]=i(G_{(i)}(x,x+e_\mu)-G_{(i)}(x+e_\mu,x))=0$: excluded.
\end{proof}
By (ii) a first-moment extraction is needed only for the one-loop bubble, so the face piece below exists only there; the local part of the vacuum polarisation (the zeroth moment of the part with lines $C^{\xi}$ only) vanishes identically on $\Omega$, faces included, by gauge invariance of the scalar functional under constant shifts of the photon field on the simply connected cube (an identity for the full kernel, not a per-ordering one), so no $D=-1$ photon piece carries a face weight.

\begin{lemma}[the face piece]\label{lem:a02:F1}
Let $\Pi^\xi(b,b')$ be the one-loop bubble with two lines $C^\xi_D\defeq(-\Delta^\Omega_D+\rho^{-2})^{-1}$ minus the seagull, and $M(x)\defeq\sum_{b'}\Pi^\xi(b,b')(x'-x)_s$ its first moment at $b=(x,\mu)$. Then $|M(x)|\le C(1+\ln_+(\rho/t(x)))e^{-\delta t(x)/\rho}+Ce^{-\delta p/\rho}$, uniformly in $5^{-j}\le\rho$.
\end{lemma}
\begin{proof}
$C^\xi_D(x,y)=\sum_{g}\operatorname{sgn}(g)C^\xi(x,gy)$ over the reflection group of $\Omega$. The term with both identity images and $x'$ summed over $5^{-j}\Z^3$ is the free first moment, which vanishes by the symmetries $C^\xi(-y)=C^\xi(y)$ and $\mu\leftrightarrow\mu'$ \cite[(3.27)--(3.28)]{Ba3}. With $|\nabla^aC^\xi(x)|\le C_a(|x|+5^{-j})^{-1-a}e^{-c|x|/\rho}$ \cite[Thm.~4.3.1, Cor.~4.3.3]{LL10}, each bubble term carries two lattice differences and $|\Pi_{\mathrm{free}}(x,x')|\le C(|x-x'|+5^{-j})^{-4}e^{-c|x-x'|/\rho}$. The free term over $x'$ outside $\Omega$ has $|x'-x|\ge t-5^{-j}$ and contributes $\le C(1+\ln_+(\rho/t))e^{-ct/\rho}$. A term with an image line: for a reflection $g$ in a face plane, $|x-gx'|\ge\max(|x-x'|,t)$; a term whose image line carries $\mathsf q$ of the two differences is $\le C|x-x'|^{-1-a}\max(|x-x'|,t)^{-1-\mathsf q}e^{-c\max/\rho}$, $a+\mathsf q=2$, and its first moment is $\le C[t^{-1-\mathsf q}\int_0^tr^{2-a}dr+\int_t^\infty r^{1-a-\mathsf q}e^{-cr/2\rho}dr]e^{-ct/2\rho}\le C[1+\ln_+(\rho/t)]e^{-ct/2\rho}$. Terms with two image lines or composed reflections decay faster; images across the faces of $\square'$ not on $F$ give $e^{-\delta p/\rho}$.
\end{proof}

\begin{lemma}[layer transfer]\label{lem:a02:L}
Let $\omega\in\lay\rho m\theta$, $\rho\le\rho'\le1$, $\delta>0$, $\delta_2\defeq\min(\delta,\delta_b)/2$. Then (a) $\sup_ye^{\delta_2t(y)/\rho'}\rho'^{-3}\sum_x5^{-3j}e^{-\delta|x-y|/\rho'}|\omega(x)|\le C\rho/\rho'$; (b) $\sum_{x\in Q}5^{-3j}|\omega(x)|\le Cs^2\rho$ for every cube $Q$ of side $s\ge\rho$.
\end{lemma}
\begin{proof}
$t(x)\ge t(y)-|x-y|$ gives $\delta|x-y|/2\rho'+\delta_bt(x)/2\rho\ge2\delta_2t(y)/\rho'$, which produces $e^{-\delta_2t(y)/\rho'}$ from half of each exponential. For the rest, $t(x)=\min_P\dist(x,P)$ over the (at most six) face planes, the profile $f(s)\defeq(1+\ln_+(\rho/s))^m(\rho/s)^\theta e^{-\delta_bs/2\rho}$ is decreasing, and for one plane the sum factorises: $\sum_{x_\parallel}5^{-2j}e^{-\delta|x_\parallel-y_\parallel|/2\rho'}\le C\rho'^2$ and $\sum_{n\ge1}5^{-j}f(n5^{-j})\le C\rho$ because $\int_0^\infty(1+\ln_+(1/u))^mu^{-\theta}e^{-\delta_bu/2}du<\infty$ for $\theta<1$. (b) is the same with the cube's cross-section.
\end{proof}

\begin{theorem}[the ordered bound with layered vertices]\label{thm:a02:OFB}
Let $G$ be a generalised graph of the cube expansion with its divergent sub-stages rewritten, every line obeying the slice bounds on $\Omega$ up to $F$, vertex functions bounded by sup norms except finitely many layered vertices $v_1,\dots,v_r$ with ultraviolet lines, and assume \DEGB: for every admitted ordering and every stage, every connected component $H$ has $D(H)\ge0$, with $D(H)=0$ only if $H$ contains a layered vertex; let $\alpha_0^B>0$ be the minimum positive degree. (i) With one vertex $x_0$ held (or none), the sum over orderings, compatible indices and positions of $|E^{(j)}(G)|$ is $\le C\eta^{\rho_0}p^{n_1}(1+k)^{p_0}\exp[-\delta\dist(\text{external})]\Lambda(x_0)$ with $\Lambda(x_0)\le C(1+\ln_+(1/t(x_0)))^{\nbar+1}t(x_0)^{-\theta_{\max}}e^{-\delta_\Lambda t(x_0)}$, $\Lambda\le C$ if $t(x_0)\ge1$, $\sum_{x_0\in Q}5^{-3j}\Lambda(x_0)\le C$ per unit cube, $\Lambda=1$ if no vertex is held. (ii) Let $0<\gamma<\min(\alpha_0^B,1-\theta_{\max})$. If every line, vertex function and layered weight obeys a pair bound $|F'-F|\le5^{-\gamma j}\rho_F^{-\gamma}\times$(the bound of $F$), a layered weight's difference lying in $\lay{\rho_v}m{\theta_v+\gamma}$, then $|E'^{(j+1)}(G)-E^{(j)}(G)|\le C5^{-\gamma j}\times$(the bound (i)).
\end{theorem}
\begin{proof}
(i) Run the multiscale induction of \cref{thm:a02:S3}(a) with one change: when the line $l(t)$ being shrunk has an endpoint composite carrying a layer weight of width $\rho_w\le\rho_{i_t}$ and its other endpoint is not held, sum the layered endpoint, splitting the retained half-exponential of $l(t)$ into three parts (the sum, the transfer of the summed endpoint's other half-exponentials, the layer transfer) and using \cref{lem:a02:L}(a) in place of the plain sum: the factor is $C\rho_{i_t}^3(\rho_w/\rho_{i_t})$, and the surviving vertex carries a layer weight of width $\rho_{i_t}$. If both endpoints are layered, sum one and keep the narrower weight; if the other endpoint is $x_0$, sum nothing new and move the weight to $x_0$. The gains of a layered composite created at step $t_1$ and summed at $t_2<\dots<t_q$ multiply to $\le\prod_{t=t_1}^{t_q-1}(\rho_{i_t}/\rho_{i_{t+1}})$, i.e.\ $+1$ in the exponent at every step in $[t_1,t_q)$ (an exact rewriting, each ratio $\le1$); if the composite is never merged into $x_0$, the last step is the volume sum, \cref{lem:a02:L}(b). At every step the exponent is $D(H_t)$ plus the number of active layered composites, which by \DEGB\ is $\ge\min(\alpha_0^B,1)>0$, except at steps where a $D=0$ component contains $x_0$ and an inactive layered composite: there the layer factor $e^{-\delta't(x_0)/\rho_*}$ at $x_0$ forces $\rho_*\ge ct(x_0)$ up to a summable tail, the number of admissible index tuples for the at most $\nbar$ such steps is $\le C(1+\log_5(1/t(x_0)))_+^{\nbar}$, the sum over the merge index gives one more logarithm and $e^{-\delta_\Lambda t(x_0)}$, and a $\theta$-weight adds $t(x_0)^{-\theta}$: this is $\Lambda(x_0)$. The infrared stage is unchanged, no infrared slice meeting a layered vertex. (ii) \cref{thm:a02:S3}(b) with (i) in place of the plain bound: the finer slices carry $\rho_i^{D}=\rho_i^{D-\gamma}\rho_i^\gamma$ with $\rho_i^\gamma\le5^{-\gamma j}$ for $i\le0$, the exponents $D-\gamma$ and $D+1-\gamma$ staying positive; the replacement telescoping weakens one element by $\gamma$ with the prefactor $5^{-\gamma j}$, a weakened layered weight being of class $\theta_v+\gamma<1$, still admissible in \cref{lem:a02:L}.
\end{proof}

\begin{theorem}[(B5$'$), photon part]\label{thm:a02:BVP}
Let a member family of the cube expansion be localised in zone~B, every photon two-point prefix component $S'$ rewritten by the zeroth- and first-moment extraction of \cite[(3.26)--(3.31)]{Ba3} with the lines $G_{j_0}=C^\xi_D+G_{j_0}(1-m^2_{j_0}-a_{j_0}P_{j_0})C^\xi_D$, and every other divergent source-free stage rewritten as in \cref{def:a02:RC}. Inside each collapsed piece write every cube line as free line plus image part and every internal sum over $\Omega$ as the sum over $5^{-j}\Z^3$ minus the sum outside, and rewrite the free piece by the free rules. Then: (a) the cube is compared with the free lattice only inside the collapsed pieces and at the two cancellation sites of the bulk structure (the free counterterm against the scalar local parts, and the reflection symmetry for the first moment), piece by piece: the zeroth moment of the $C^\xi_D$-part vanishes on both sides; the residual $\Delta c^{S'}$; the face piece $M$; and the remaining differences. (b) Each difference piece is a layered vertex at stage $j_0(S')$: $M\in\lay{\rho_{j_0}}10$ (\cref{lem:a02:F1}), one leg differentiated, of the degree of the bulk first-moment vertex; $\Delta c^{S'}$, with $\sup\le C\rho_{j_0}^{-1}e^{-\delta t/\rho_{j_0}}$ (one line replaced by its image part, $\sum_w5^{-3j}|\nabla G^{\mathrm{img}}_{j_0}(x,w)|\le C\rho_{j_0}e^{-\delta't/\rho_{j_0}}$ since $|x-gw|\ge\max(|x-w|,t)$), of class $\lay{\rho_{j_0}}00$ and degree $D(S')$; the remaining pieces, ultraviolet-finite at $j_0$ with an image line, of class $\lay{\rho_{j_0}}00$ by the per-scale sum of \cite[(3.17)--(3.18)]{K2} in the $D>0$ case, with degrees $\ge D(S')+1+\alpha$ (Taylor remainders) and $\ge1$ (block members, \cref{lem:a02:BLK}). No piece has weight $t^{-1}$. (c) \DEGB\ holds for every difference graph (\cref{lem:a02:CENpp}(iii), the rewriting never lowering the degree of a containing stage), so \cref{thm:a02:OFB} applies: the pair bound with $5^{-\gamma j}$, the face piece's difference being of class $\lay{\rho_{j_0}}1\gamma$ times $5^{-\gamma j}\rho_{j_0}^{-\gamma}$ by the proof of \cref{lem:a02:F1} with one line replaced by its difference. (d) In a host rooted at an interior source vertex ($t(x_0)\ge\eta^{-1}/2$), $\Lambda(x_0)\le C$: the zone-B photon pieces cost a constant. For the source-free field-dependent classes, root-summed, $\sum_{x_0}5^{-3j}\Lambda(x_0)\le C$ per unit face cell (\cref{lem:a02:L}(b)), the pair difference carries $5^{-\gamma j}$, and the zone-B total is $\le C|\partial T^{(j)}|\le C\eta|T^{(j)}|$ relative to the bulk form.
\end{theorem}
\begin{proof}
(a) The image representations are exact and the rewriting rules are identities on each side; the difference of two rewritten pieces is the sum of the differences of their parts. (b) \cref{lem:a02:F1}, the one-loop computation with one image line, and the per-scale sum with \cref{lem:a02:CENpp}(ii) (no $D=-1$ piece other than the one-loop class, whose $C^\xi_D$-zeroth moment is exactly zero). (c) \cref{lem:a02:CENpp}(iii) and the degree bookkeeping of the pieces. (d) \cref{thm:a02:OFB}(i), (ii) and \cref{lem:a02:L}(b).
\end{proof}

\subsection{Assembly}\label{ssec:a02:assembly}

\begin{proof}[Proof of \cref{prop:k01:Kpos}]
Fix $J\ge J_0$, $f$, $x$, $s\in\{u_x,0\}$, and $k\ge k_0$, $J_*\le j\le k-J_*$ with $\eta\le\eta_1$, where $k_0$, $\eta_1$ are such that $j\le k-J_0$, $\eta\le\min(\eta_0,\delta_*)$ (\cref{lem:a02:D}, \cref{lem:a02:Z}) and $J_*\ge1$. On $\supp\chi_j$ write $S^{(j)}_{N,k}-S'^{(j+1)}_{N,k+1}$ by \eqref{eq:a02:390} for both members; the constants $\ln N_j$, $\ln N_j(0)$, $\ln N'_{j+1}$, $\ln N'_{j+1}(0)$ and the field- and source-free parts of $P^{(j),1}$, $P'^{(j+1),1}$ are in $v_{k,j}$, $v_{k+1,j+1}$ and drop out. The remaining difference is the sum of (B1)--(B5$'$), with (B6) inside (B3)--(B5$'$):
\begin{itemize}[nosep]
\item (B1) $\le C5^{-2j}p^2(\lmu\eta)^{-2}\eta^{-3}|\CJ|$ by \cref{prop:a02:B1};
\item (B2) $\le C_\varsigma p^{c+2}\lambda^{-1/2}(\lmu\eta)^{-1}\eta^{-7/2}(5^{-\gamma'j}+\eta^\varsigma)|\CJ|$ by \cref{cor:a02:B2} (its zero-background part also by \cref{prop:a02:S}(c));
\item (B3) $\le CNe^2p^{c+2}(\lmu\eta)^{-1}\eta^{-2}(5^{-\gamma'j}+5^{-j}+\eta^\varsigma)|\CJ|$ by \cref{cor:a02:B3};
\item (B4) and (B5): the curvature classes and the source-free field-dependent classes; their zone-I localisations by \cref{thm:a02:B4} (through \cref{lem:a02:I}, \cref{lem:a02:SBP}, \cref{thm:a02:T}, \cref{thm:a02:S3}), their zone-B tails (the $R>p$ branch of a host rooted at the interior source vertex, \cref{lem:a02:Z}(ii)) by \cref{thm:a02:Bsc}(a) and \cref{thm:a02:BVP}(d) at $\Lambda\le C$ with the factor $e^{-\delta R/2}$, and the position sum by \cref{lem:a02:SC}: together $\le C[5^{-\gamma'j}\eta^{-\beta'}+\eta^{\sigma'}]|\CJ|$;
\item (B5$'$): the source-free field-dependent classes localised in zone~B, with the rate, by \cref{thm:a02:Bsc}(b),(c) and \cref{thm:a02:BVP}(c),(d) (mixed members with scalar remainders and photon face pieces included): $\le C5^{-\gamma'j}\eta^{-\beta'}|\CJ|$;
\item (B6): the members with block vertices are bounded inside each of the above through \cref{lem:a02:BLK} and \cref{thm:a02:S3}, which supply the degree and the ordered bound in the split count; their pair bound is \cref{thm:a02:S3}(b).
\end{itemize}
Every piece is of the form $C(5^{-\gamma j}\eta^{-\beta}+\eta^\sigma)$ for some $\gamma\in(0,1)$, $\beta,\sigma>0$ (the factor $|\CJ|$ and the powers of $p=p(\eta)$, $(1+k)$ are absorbed on the range $j\le k$, $\eta=5^{j-k}$, into $\eta^{-\beta}$ by an arbitrarily small loss in $\beta$, and $5^{-\gamma'j}$ with $\gamma'\ge1$ is weakened to $\gamma<1$). Take the smallest $\gamma$, the largest $\beta$ and the smallest $\sigma$ among the pieces, and the sum of the constants as $C_3$. The constants depend on the point of King's range, $J$ (through $|\CJ|$ and $J_0$), $f$ (through $\max\norm{f_{\mu\nu}}_{C^{2,\alpha}}$ and $\beta_f$) and $x$ (through $|x|$ in the source vertex), and on nothing else. This is \eqref{eq:k01:Kpos}.
\end{proof}
